% Transcription — fonds Alexandre Grothendieck, shelfmark 1, batch 2 (pages 21-40).
% Established from the facsimile archives/batches/1/batch-02.pdf (a mirror of
% https://grothendieck.umontpellier.fr/1.pdf), per the skill
% transcribe-grothendieck.
%
% Pass: Fable 5.1 (claude-fable-5-1), 2026-09-05 — first pass, unchecked against
% the pages by a human.
%
% Nineteen of the twenty pages are transcribed. Page 36 is blank and absent.
%
% \page{} numbers the position in the folder, as batch 1 did: the archivists'
% pencil numbers run two ahead from sheet 3 of the folder on, so page 21 here
% is pencilled 23 and page 40 is pencilled 42. The note at page 3 of batch 1
% records this once for the folder.
%
% The batch falls into three runs, all in the same fast blue-black hand, and
% the prose between the formulas is largely lost — the apparatus says so
% rather than inventing. The statements come through; the connective
% sentences are \uncertain{} word by word or \ill{}.
%
%   Pages 21-33 — the notes continue from batch 1's section 4. Page 21 ends
%     the proof of Proposition 1 (when K(M) of a maximal commutative
%     subalgebra falls in the ideal a) and states Corollaire 1 (the trace
%     restricted to C_0(M) is a measure μ, T(f) = ∫ f dμ). Page 22: Cor. 2,
%     the spectral measure μ_A of a normal A ∈ b, (1) ⟨φ, μ_A⟩ = T(φ(A)),
%     (3) ∫ z^n dμ_A = T(A^n), the spectral rearrangement φ_A of A ≥ 0 and
%     (4) Φ_A, Ψ_A, Δ_A. Page 23: (5)-(6), Δ(A) = Δ_A(ω) with ω = T(1);
%     then, cancelled by diagonals and taken up on 24-25, the decomposition
%     of the algebra along the mass points of M into type I factors and the
%     formula T(f) = ∫_0^1 T(f(s)) ds. Pages 26-27: Prop 6 (φ_A^-(t) as an
%     inf over projectors of trace < t of ‖(1-E)A(1-E)‖), its four
%     corollaries (φ_{|BA|} ≤ ‖B‖φ_{|A|}, monotonicity, semicontinuity),
%     Prop 7 (φ_{|A|} = φ_{|A*|} via T((AA*)^p) = T((A*A)^p)). Then his
%     « N° 6 » Compléments sur le déterminant: (1) Δ(A) = exp ∫_0^{T(1)}
%     log φ_{|A|}, the three cases |λ| < 1, > 1, = 1 for A = λ1 + B,
%     Prop 8 (Δ(A) = Δ(A*), Δ(AB) = Δ(A)Δ(B) on 1 + a) with a first proof
%     cancelled on page 29, the relative determinants Δ^{X,Y} through
%     partial isometries on page 30, Prop 9 (Δ_A(t) = sup_{Tr E ≤ t}
%     Δ^E(EAE)) and Théorème 1 (Δ_{AB}(t) ≤ Δ_A(t)Δ_B(t)) on pages 30-32.
%     Page 32 opens « Inégalités fondamentales » — the plan's N° 6, « Les
%     inégalités pour AB »: Th. 2 (W(AB) ≤ W(φ_{|A|}φ_{|B|}) for W an
%     increasing Weyl function), the boxed (4) ∫ φ_{|AB|} ≤ ∫ φ_{|A|}φ_{|B|},
%     Hölder for operators (6) ‖AB‖_r ≤ ‖A‖_p ‖B‖_q, the duality (7)
%     Φ_{|A|}(t) = sup |Tr AB| over ‖B‖ ≤ 1, Tr σ(B) ≤ t, and the
%     subadditivity (8) Φ_{|A+B|} ≤ Φ_{|A|} + Φ_{|B|}.
%   Pages 34-35 — a second mouture of the same material on smaller sheets,
%     opening with the word « Astuce »: Prop 1 (φ_A(t+0) as the same inf),
%     Cor. 1-5, Prop 2 (φ_{|A|} = φ_{|A*|}), then a Proposition giving
%     Δ_t(A) and Δ_t(1+|A|) as sups over projectors. Transcribed apart,
%     not merged with the first.
%   Pages 37-40 — four yellowed sheets of another set, boxed in red pencil:
%     Préliminaires on compact operators u : E → F between Hilbert spaces,
%     polar decomposition u = U|u|, |u| = Vu, the singular values ρ_n(u) =
%     λ_n(|u|) and their minimax characterisation, ρ_n(uv) ≤ ‖u‖ρ_n(v), Ky
%     Fan's inequalities ρ_{m+n-1}(AB) ≤ ρ_m(A)ρ_n(B); « II »
%     Caractérisation des op. de Fredholm, ‖u‖_1 = Σ ρ_i(u), Σ_1^n ρ_i(u) =
%     sup |Tr v_n u| over rank ≤ n, ‖u+v‖_1 ≤ ‖u‖_1 + ‖v‖_1; and the
%     definition of L^p(E,F) by ‖u‖_p = (Σ ρ_i(u)^p)^{1/p}. Page 39 is
%     largely cancelled by zigzags; page 40 carries three lines.
%
% Cancellations: pages 23, 24, 25, 28, 29, 31 and 39 carry blocks struck
% through by diagonals or loops; each is transcribed as \struck{} with one
% \note{} for the block. The red pencil frames on pages 37-40 are recorded
% once per formula. A « Rq » and a « ? » in red pencil stand in the margins
% of pages 37-38.
%
% Notation is that of batch 1: \mathcal{A} the von Neumann algebra, T its
% trace, \underline{a} the ideal of definition, \underline{b} its norm
% closure, \underline{C} a commutative subalgebra, \mathcal{P} the
% projectors, \varphi_A the spectral rearrangement, \Delta the determinant.
% He writes Δ_A(t) up to page 33 and Δ_t(A) from page 35; Δ_E and Δ^E both
% occur for the relative determinant, the superscript replacing a struck
% subscript on pages 29-30. Each is transcribed as written.
%
% The dating is the inventory's own for the shelfmark. Its group,
% « MATHÉMATIQUES / RECHERCHE / Travaux », carries no date.
\documentclass[11pt,a4paper]{article}
% Le préambule commun à toutes les transcriptions du fonds Grothendieck.
%
% Il tient en une idée : **l'incertitude se compose, elle ne se résout pas.**
% Une transcription qui lisse un mot douteux a détruit la seule chose pour
% laquelle on la faisait — un lecteur doit pouvoir savoir, à chaque endroit,
% ce qui était sur le feuillet et ce qui a été ajouté. Les six macros qui
% suivent sont donc l'appareil critique, et non de la décoration.
%
% Les mêmes commandes sont reconnues par `scripts/render.mjs`, qui produit la
% vue de lecture du site. Une macro ajoutée ici doit l'être aussi là-bas,
% faute de quoi le rendu échouera bruyamment — ce qui est voulu : un rendu
% silencieusement faux est pire qu'un rendu absent.

\ProvidesPackage{grothendieck}[2026/08/08 Transcriptions du fonds Grothendieck]

\usepackage{amsmath,amssymb,amsthm}
\usepackage{mathrsfs}
\usepackage{stmaryrd}
% Les diagrammes commutatifs sont partout dans ce fonds : c'est la langue
% même de la géométrie algébrique de Grothendieck, et une transcription qui
% les rendrait en prose ne transcrirait rien.
\usepackage{tikz-cd}
% Français par défaut — c'est la langue des feuillets — mais l'anglais est
% chargé aussi : la traduction et le résumé sont des documents anglais, et la
% typographie française (espaces avant les deux-points) y serait une faute.
% Ces documents basculent par \selectlanguage{english} en tête de corps.
\usepackage[english,french]{babel}
\usepackage{fontspec}
\usepackage{geometry}
\geometry{a4paper,margin=25mm,marginparwidth=28mm}
\usepackage{marginnote}
\usepackage{xcolor}
% ulem, not soul. `soul`'s \st and \ul are built on a character-by-character
% scan that XeLaTeX's font machinery breaks: the document fails with an
% undefined control sequence rather than degrading. ulem does the same two jobs
% and is Unicode-engine safe. `normalem` stops it hijacking \emph, which the
% transcriptions use for Grothendieck's own emphasis.
\usepackage[normalem]{ulem}
\usepackage{hyperref}
\hypersetup{colorlinks=true,linkcolor=black,urlcolor=black,pdfborder={0 0 0}}

% --- Métadonnées du lot -----------------------------------------------------
% Reprises telles quelles par le script de rendu, qui les affiche en tête de
% la vue de lecture. Elles fixent ce que le fichier prétend couvrir : sans
% elles, une transcription flotte sans point d'attache dans un fonds de seize
% mille feuillets.

\newcommand{\folder}[1]{\gdef\grothFolder{#1}}
\newcommand{\batch}[1]{\gdef\grothBatch{#1}}
\newcommand{\leaves}[2]{\gdef\grothFirst{#1}\gdef\grothLast{#2}}
\newcommand{\foldertitle}[1]{\gdef\grothFoldertitle{#1}}
\gdef\grothFolder{?}\gdef\grothBatch{?}\gdef\grothFirst{?}\gdef\grothLast{?}\gdef\grothFoldertitle{}

% --- Le résumé --------------------------------------------------------------
%
% Il ouvre la lecture modernisée, sous le titre « Résumé », et il est composé
% en petit. Ce n'est pas un ornement : c'est le seul passage du document qui
% ne soit pas des mathématiques, et un lecteur doit pouvoir le reconnaître
% avant de l'avoir lu — puis le sauter s'il connaît déjà le sujet, ou s'y
% arrêter s'il ne le connaît pas. Le filet qui le suit marque la reprise du
% corps.

\newenvironment{resume}
  {\small\setlength{\parskip}{0.45em}}
  {\par\medskip\hrule\medskip}

% --- Mots-clés ---------------------------------------------------------------
%
% En anglais, en clôture du résumé de la lecture modernisée : le vocabulaire
% moderne sous lequel chercher ce que les feuillets construisent. Le manifeste
% du site (`npm run manifest`) les extrait pour étiqueter la cote entière —
% cette ligne est la source unique des tags, il n'y a pas de fichier à côté.
\newcommand{\keywords}[1]{\par\smallskip\noindent
  {\footnotesize\textsc{Keywords} --- \textit{#1}}\par}

% --- L'appareil critique ----------------------------------------------------

% Le numéro de feuillet, en marge. C'est l'ancre qui permet de régler un
% désaccord sur une formule en regardant *un* feuillet plutôt qu'une chemise
% de six cents. Le site s'en sert aussi pour tourner les pages du fac-similé
% quand on fait défiler la transcription.
\newcommand{\leaf}[1]{%
  \marginnote{\footnotesize\color{gray!70}\textsf{#1}}%
  \label{leaf:#1}\ignorespaces}

% Illisible : on ne devine pas. Un mot inventé qui se lit comme les autres est
% le pire résultat possible.
\newcommand{\ill}{\textcolor{red!60!black}{[\,\ldots\,]}}

% Lecture douteuse : le mot est donné, et signalé comme incertain.
\newcommand{\uncertain}[1]{\uline{#1}}

% Ajout de l'éditeur — une lettre manquante, un mot sous-entendu.
\newcommand{\add}[1]{\textcolor{blue!50!black}{[#1]}}

% Biffé par Grothendieck lui-même. Ce qu'il a rayé fait partie du document :
% une rature dit souvent où la pensée a changé de direction.
\newcommand{\struck}[1]{\sout{#1}}

% Note du transcripteur, sur l'état matériel du feuillet.
\newcommand{\note}[1]{\par\smallskip{\small\color{gray!60!black}\textit{[#1]}}\par\smallskip}

% Note marginale de Grothendieck, distincte du corps du texte.
\newcommand{\marginal}[1]{\par\smallskip\noindent\hspace{1em}%
  {\small\color{gray!60!black}#1}\par\smallskip}

% --- Titre ------------------------------------------------------------------

\AtBeginDocument{%
  \begin{center}
    {\large\bfseries Fonds Alexandre Grothendieck --- cote n\textsuperscript{o}~\grothFolder}\\[2pt]
    {\itshape\grothFoldertitle}\\[4pt]
    {\small feuillets \grothFirst--\grothLast\ (lot \grothBatch)}\\[2pt]
    {\footnotesize\color{gray!60!black}Transcription établie d'après le fac-similé numérisé
     par l'Université de Montpellier.\\
     Les lectures incertaines sont soulignées, les illisibles notées [\,\ldots\,],
     les ajouts entre crochets.}
  \end{center}
  \vspace{1em}}

\endinput


\folder{1}
\batch{2}
\pages{21}{40}
\dating{1953}
\watermark{Édition de démonstration}
\foldertitle{Analyse fonctionnelle : notes manuscrites (s.d.), lettre (1953)}

\begin{document}

\section{Mesures spectrales et réarrangements spectraux d'opérateurs (suite)}

\note{Titre de l'éditeur, repris du titre de sa main en tête de la page 18 (lot 1). La démonstration de la proposition 1, commencée page 20, s'achève page 21.}

\page{21}
\note{Le lot précédent s'arrête au milieu de la démonstration de la proposition 1 du N° 4 ; elle se poursuit ici, en tête d'une page dont les cinq premières lignes sont biffées.}

\struck{\ill{} $U_{\varphi_K}$ \ill{} dans $\underline{b}$, \ill{}} Soit $f$ \uncertain{fonction} \uncertain{hermitienne} \uncertain{bornée} sur $M$ \ill{} : i.e. \struck{l'opérateur $U_f$ est dans $\underline{b}$ : $f$ \ill{} \uncertain{majorée} \ill{} $|U_f|$ \ill{} $U_{|f|}$ \ill{} $g \in C_0(M)$,} \uncertain{Or} \uncertain{si} $f$ \ill{} \uncertain{majorée} \ill{} \uncertain{par} \uncertain{un} $g \in C_0(M)$, \uncertain{d'où} $0 \leq |U_f| \leq U_g \in \underline{b}$, \uncertain{donc} $U_f \in \underline{b}$, i.e. $U_f$ \uncertain{est} \ill{} \ill{} $U_h$, $h \in C_0(M)$, \uncertain{un} $E$ \uncertain{projecteur}, \uncertain{est} \ill{} \uncertain{si} $f$ \uncertain{est} \ill{} \struck{\ill{} $\varphi_K$, \ill{}} \uncertain{borné} $\varphi_K$, $K$ \uncertain{compact} $\subset M$. \uncertain{Donc} il \uncertain{existe} alors un \uncertain{compact} \ill{} $K'$ \uncertain{tel} \uncertain{que} $U_{\varphi_K} = U_{\varphi_{K'}}$, \ill{} $K \subset K'$. \ill{} $K$ \ill{} \ill{} \uncertain{et} \uncertain{soit} $E = U_{\varphi_K}$. \uncertain{Tout} \uncertain{opérateur} $A$ dans $\mathcal{A}$ qui \uncertain{est} \uncertain{permutable} à $\underline{C}$, et tel que $EAE = A$, est \uncertain{dans} $\underline{C}$, et $=$ \uncertain{dans} $C(K)$. \uncertain{Donc} $\underline{C}(K)$ \uncertain{est} \ill{} \uncertain{et} \uncertain{contient} \uncertain{l'algèbre} \uncertain{formée} \uncertain{des} \uncertain{opérateurs} \uncertain{permutables}, \uncertain{algèbre} qui est \struck{\ill{}} \uncertain{plus} \uncertain{fine} \uncertain{dans} $\mathcal{A}$ \uncertain{et} \uncertain{Neumann}. \note{de « Tout opérateur » à « Neumann », les mots pris pour \emph{permutable}, \emph{contient}, \emph{formée} et \emph{plus fine} ne sont que des candidats ; l'argument est que le commutant de $\underline{C}$ dans $E\mathcal{A}E$ se réduit à $C(K)$}

\emph{Corollaire 1} \uncertain{Donc} \uncertain{si} \ill{} $\underline{C}$ \uncertain{est} \uncertain{l'algèbre} \uncertain{de} \uncertain{définition} \uncertain{d'une} \uncertain{trace} \note{en interligne, au-dessus de « trace » : « \uncertain{semi-finie} »} $T$, $\underline{C}$ \uncertain{une} \uncertain{sous-algèbre} \uncertain{commutative} \uncertain{maximale} \ill{} \uncertain{de} $\underline{b}$ \uncertain{et} \ill{} \uncertain{de} $\underline{a}$, $\underline{C} = C_0(M)$, alors \ill{} \uncertain{une} \uncertain{mesure} \uncertain{positive} $\mu$ \note{le $\mu$ est surchargé} telle que \uncertain{si} $f \in C_0(M)$, \ill{} \uncertain{soit} $f \in \underline{a}$ \uncertain{si} \uncertain{et} \uncertain{seulement} \uncertain{si} \ill{} $f$ \ill{} \uncertain{soit} \uncertain{intégrable} \ill{}. Alors $T(f) = \int f\, d\mu$. \uncertain{Si} $T$ \uncertain{est}
\note{une accolade en marge gauche embrasse le corollaire jusqu'au bas de la page ; la phrase se poursuit page 22}

\page{22}
\emph{Corollaire 2} Soit $A$ un \uncertain{opérateur} \uncertain{normal} \uncertain{de} $\underline{b}$, $f$ une \uncertain{fonction} \uncertain{continue} \uncertain{sur} $\underline{\mathbf{C}}$ \uncertain{bornée}, \ill{} \uncertain{compacts}, nulle \uncertain{au} voisinage de $0$. Alors $f(A) \in \underline{a}$.

Soit \uncertain{maintenant} $A$ un \uncertain{opérateur} \uncertain{normal} de $\underline{b}$. $A$ \uncertain{est} \uncertain{contenu} \uncertain{dans} \uncertain{une} \uncertain{sous-algèbre} \uncertain{commutative} \note{en interligne, cerné : « \ill{} \uncertain{involutive} »} \uncertain{maximale} \ill{} \uncertain{de} $\underline{b}$, \uncertain{qui} \uncertain{s'identifie} \uncertain{à} \uncertain{un} $C_0(M)$, \uncertain{sur} $M$ \uncertain{il} \uncertain{y} \uncertain{a} \uncertain{une} \uncertain{mesure} \struck{\uncertain{$\mu_A$}} \uncertain{La} \uncertain{mesure} \uncertain{spectrale} $\mu_f$ \uncertain{de} $f$ \uncertain{existe}, \uncertain{elle} \uncertain{ne} \uncertain{dépend} \uncertain{pas} \uncertain{du} \uncertain{choix} \uncertain{de} $\underline{C}$, \uncertain{car} \uncertain{on} \uncertain{a} \uncertain{pour} \uncertain{tout} $\varphi \in K(\mathbf{C}^{*})$ :
\[
  \langle \varphi, \mu_f\rangle = \int_M \varphi(f) = T(\varphi(A))\,.
\]
\uncertain{On} \uncertain{l'appelle} \emph{\uncertain{mesure} \uncertain{spectrale} de $A$}, \uncertain{et} \uncertain{on} \uncertain{la} \uncertain{note} $\mu_A$ :
\[
  \text{(1)}\qquad \langle \varphi, \mu_A\rangle = T(\varphi(A)) \qquad \varphi \in K(\mathbf{C}^{*})
\]
\struck{\uncertain{si} $\varphi$ \uncertain{est} \uncertain{finie}, \uncertain{la} \ill{}} \struck{\ill{} ① \ill{} $\mu_A = \sigma(A)$ (\uncertain{spectre} de $A$).} \note{la ligne encadrée porte un « ① » cerclé ; le premier membre, illisible, est le support ou la norme de $\mu_A$ ; une longue accolade en marge gauche embrasse les cinq lignes qui suivent} \uncertain{Ce} \ill{} \uncertain{résulte} \ill{} \uncertain{une} \uncertain{égalité}, \uncertain{quand} $\varphi$ \uncertain{est} \uncertain{finie}. \uncertain{Le} \uncertain{volume} \uncertain{résulte} \uncertain{que} $A \in \underline{a} \Longleftrightarrow$ \ill{} $2$ \ill{} \struck{\ill{} $\mu_A$, \uncertain{Alors} \ill{} il \uncertain{existe} $\int z^n\, d\mu_A$}
\[
  \text{(3)}\qquad \int z^n\, d\mu_A(z) = T(A^n) \qquad (n \geq 1)
\]
\note{le numéro est un « ③ » cerclé ; aucun (2) n'apparaît}

\marginal{en marge gauche, sideways, sur neuf lignes : \uncertain{le} \ill{} \ill{}, \uncertain{identique} \ill{}, \uncertain{la} \uncertain{mesure} \uncertain{de} \ill{} \ill{}, \ill{} \uncertain{et} \ill{}, \uncertain{sur} $B$, \ill{} \uncertain{Corollaire} \ill{} \uncertain{i.e.} \ill{}}

\uncertain{Supposons} \uncertain{maintenant} $A \geq 0$. Alors \struck{$\mu_A$} $f$ \ill{} \uncertain{est} \uncertain{décroissante} \ill{}, \uncertain{qui} \ill{} \uncertain{ne} \uncertain{dépend} \uncertain{que} de $A$ (\uncertain{car} $\mu_A$ \uncertain{ne} \uncertain{dépend} \uncertain{que} \uncertain{de} $A$), \ill{} \uncertain{soit} $\varphi_A$, \uncertain{et} \uncertain{appelé} \emph{\uncertain{réarrangement} \uncertain{spectral} de $A \in \underline{b}^{+}$} (\uncertain{c'est} \uncertain{une} \uncertain{fonction} \uncertain{décroissante} \uncertain{sur} $\mathbf{R}^{+}_{*}$, $\|\varphi_A\| = \|A\|$). \uncertain{On} \uncertain{voit} \uncertain{aussitôt} \uncertain{que} $\Phi_A = \Phi_f$, $\Psi_A = \Psi_f$ (\uncertain{donc} \ill{} \ill{})
\[
  \text{(4)}\qquad \left\{
  \begin{array}{l}
    \Phi_A(t) = \displaystyle\int_0^t \varphi_A(s)\, ds \\[1ex]
    \Psi_A(t) = \displaystyle\int_0^t \log\varphi_A(s)\, ds \\[1ex]
    \Delta_A = \Delta_f = \exp\Psi_f \qquad \Delta_A(t) = \exp\displaystyle\int_0^t \log\varphi_{|A|}(s)\, ds
  \end{array}\right.
\]
\note{à gauche de (4), une première version « $\Phi_A = \Phi_f$, $\Psi_A = \Psi_f$ » est biffée d'un gribouillis ; la dernière ligne, écrite sous l'accolade, est soulignée d'un trait qui court jusqu'au bord droit}

\page{23}
\begin{align*}
  \text{(5)}\qquad & \Delta_A(t) = \Delta_{|A|}(t) = \Delta_{\varphi_{|A|}}(t) = \exp\int_0^t \log\varphi_{|A|}(s)\, ds \\
  \text{(6)}\qquad & \Delta(A) = \struck{\exp\Psi_A} \; \Delta_A(\omega) \qquad (\text{si } \omega = T(1) < +\infty)
\end{align*}
\note{le membre biffé de (6) est un gribouillis dont on lit un « exp » ; une longue accolade en marge gauche embrasse (5), (6) et les quatre lignes suivantes}

\uncertain{Par} \uncertain{suite}, $\Delta_A(t)$ \uncertain{est} \uncertain{une} \uncertain{fonction} \uncertain{continue} \uncertain{de} $t$ : \uncertain{elle} \ill{}, \uncertain{si} $\varphi_{|A|}(s) < \infty$. \struck{$\Delta(A)$ \uncertain{est} \uncertain{fini} \ill{} \uncertain{si}} $T(1) < +\infty$ (i.e. \uncertain{si} $T$ \uncertain{est} \uncertain{une} \uncertain{trace} \uncertain{finie}), \uncertain{alors} \ill{} \ill{} \uncertain{fini}, \ill{} $\uncertain{0} \leq \Delta(A) < +\infty$. \note{le premier membre de la dernière inégalité se lit aussi « $1$ »}

\marginal{en marge gauche, sideways, à hauteur de (5)-(6) : \uncertain{Définition} \ill{} $\displaystyle\int_1^{\infty}$ \ill{} \uncertain{$\varphi_{|A|}$} \ill{}, \uncertain{(5)} $= C_0$ \ill{}, \uncertain{$U_{\varphi}$} \ill{}, \ill{} \uncertain{$W(A)$} \ill{} « $\Delta_A$ » \ill{} \uncertain{définit} \ill{}, \ill{} $\varphi$ \ill{} ; puis, plus bas et sideways aussi : \ill{} $\displaystyle\int_0^{\infty}$ \ill{} $\mathbf{R}^{+}_{*}$ \ill{}}

\uncertain{Supposons} \uncertain{que} $\underline{C}$ \uncertain{soit} \uncertain{une} \uncertain{sous-algèbre} \uncertain{commutative} \uncertain{involutive} \uncertain{maximale} \uncertain{dans} $\underline{b}$, $\underline{C} = C_0(M)$, \uncertain{et} \uncertain{la} \uncertain{mesure} \uncertain{sur} $M$ \uncertain{définie} \uncertain{par} \uncertain{le} \uncertain{corollaire}~1. \uncertain{Si} $t \in M$ \uncertain{est} \uncertain{un} \uncertain{point} \uncertain{de} $M$ \uncertain{de} \uncertain{mesure} \uncertain{non} \uncertain{nulle}, \ill{} \uncertain{soit} \ill{} \uncertain{borné} \ill{} 2 \ill{} $\underline{C}$) \uncertain{par} $E = U_{\varphi_{\{t\}}}$ \note{le $E$ et l'indice sont surchargés} \ill{} \uncertain{un} \uncertain{projecteur} \uncertain{minimal} \ill{} \uncertain{le} \struck{\ill{}} \uncertain{définit} \ill{} \uncertain{soit} $E$ \uncertain{est} \uncertain{un} \uncertain{projecteur} \uncertain{minimal} \uncertain{dans} $\mathcal{A}$ \struck{$= \sigma(M)$ \ill{} \uncertain{algèbre} \ill{}} \ill{} \uncertain{de} $\underline{b}$, \ill{} \uncertain{contient} $E$, \uncertain{alors} \struck{\ill{}}

\struck{\ill{} \uncertain{caractéristique} \uncertain{d'un} \uncertain{point} \ill{} \ill{} \uncertain{de} $M$, \ill{} : \uncertain{un} \ill{} $\geq 0$. \uncertain{On} \uncertain{voit} \ill{} \ill{} \ill{} $E$, i.e. \ill{} \ill{} $T$ \ill{} \ill{} (\uncertain{i.e.} $E$ \uncertain{est} \uncertain{un} \ill{}) \ill{} \ill{} \uncertain{réarrangement} \ill{} \uncertain{compact} \ill{} $\underline{b}$ \ill{} \uncertain{de} $\mathcal{A}$ \ill{}}
\note{les douze dernières lignes de la page sont annulées par trois longues diagonales et cernées de traits ; une seule ligne, à mi-hauteur, échappe aux diagonales : « \uncertain{Supposons} \ill{} $\geq 0$. \uncertain{On} \uncertain{voit} \uncertain{que} \ill{} », qui reste de toute façon illisible. La question reprise page 24 est celle du projecteur minimal $E$ et de la mesure du point $t$.}

\page{24}
\note{Page entièrement en prose rapide ; seules les formules se lisent avec sûreté. La moitié inférieure, de « Soit alors » au bas de la page, est annulée par trois longues diagonales.}

\ill{} $t \in \Omega$ \uncertain{tel} \uncertain{que} $U_{\varphi_{\{t\}}} E = E\, U_{\varphi_{\{t\}}} = E$, \ill{} \note{le $\Omega$ n'est pas corrigé en $M$ ici} $\mathcal{A}$ \uncertain{est} \uncertain{alors} \struck{\ill{}} \uncertain{somme} \uncertain{directe} \uncertain{de} \uncertain{deux} \uncertain{algèbres} $\mathcal{A}_0$ et $\mathcal{A}_1$, \uncertain{où} $\mathcal{A}_1$ \uncertain{est} \uncertain{un} \uncertain{facteur} \ill{}, \uncertain{et} \ill{} \uncertain{les} \uncertain{projecteurs} \uncertain{minimaux} \ill{} $E$. \uncertain{Le} \struck{\ill{}} \ill{} \uncertain{les} \uncertain{parties} \uncertain{de} $\mathcal{A}_0$ et $\mathcal{A}_1$, \ill{} \uncertain{des} \uncertain{projecteurs} \uncertain{de} $\mathcal{A}_1$ \note{deux lignes cernées d'un long trait} \ill{} \uncertain{les} $\mathcal{A}_i$ \uncertain{sont} \uncertain{de} \uncertain{type}~I, \uncertain{c'est-à-dire} \ill{} \uncertain{isomorphes} \ill{} : \uncertain{un} $L(\mathcal{H})$, $T$ \uncertain{induit} \ill{} \struck{\ill{}} \uncertain{par} \ill{} \uncertain{la} \uncertain{trace} \uncertain{usuelle} \uncertain{dans} $\mathcal{A}_1$ (\uncertain{à} \uncertain{un} \uncertain{facteur} \uncertain{constant} \uncertain{près}), \uncertain{et} $T(E) < +\infty$, \struck{$E$ \ill{}} \ill{} $EF\underline{a} \subset \underline{b}$. \struck{\uncertain{Alors} \ill{}, \uncertain{si} $E$ \ill{}} \ill{} \uncertain{Soit} \uncertain{alors} $\underline{b} = C_0(M)$ \ill{} \uncertain{qui} \uncertain{contient} \ill{} $\underline{b}$, \ill{} \ill{} $E$, \ill{}, \uncertain{i.e.} \uncertain{la} \uncertain{définition} : \uncertain{les} \uncertain{parties} \ill{} \uncertain{de} \ill{} $M$, \uncertain{qui} \ill{} \uncertain{compact} : 2 \ill{}. \ill{} \uncertain{qui} \uncertain{sont} \uncertain{minimaux} \ill{} (\uncertain{est} \ill{} $E$ \ill{}) \ill{}, \uncertain{si} \uncertain{et} \uncertain{seulement} \uncertain{si} $m(\{t\}) > 0$. \note{le $m$ est un trait ; on lit « $\{t\}$ » en exposant d'un signe surchargé} \ill{} \uncertain{soit} \ill{}

\struck{\ill{} \uncertain{comme} $\underline{C} = C_0(M)$, \ill{} \uncertain{Alors} \ill{}, \ill{} \uncertain{soit} $M$ \ill{} \uncertain{un} \uncertain{point}, \ill{} $\mathcal{A}$ \uncertain{est} \ill{} \uncertain{par} \ill{} \uncertain{Donc}, \ill{}, \ill{} \uncertain{Soit} $I$ \uncertain{l'ensemble} \uncertain{des} \ill{} $t \in M$ \uncertain{de} \ill{} \uncertain{pour} \uncertain{tout} $t$, \uncertain{contenu} \ill{} \uncertain{de} $\mathcal{A}$ \ill{} \uncertain{algèbre} $\mathcal{A}_i = U_{\{t\}}\mathcal{A}\, U_{\{t\}}$ \uncertain{de} $\mathcal{A}$ \ill{}, \uncertain{facteur}) \uncertain{de} \uncertain{type}~I, \uncertain{i.e.} \ill{} \uncertain{et} \ill{} \ill{}. \uncertain{On} \uncertain{voit} \uncertain{donc} \ill{} $\mathcal{A}$ \ill{}}
\note{le bloc annulé est repris, dans le même ordre d'idées, page 25 : décomposition de $\mathcal{A}$ suivant les points de masse de $M$, chacun donnant un facteur de type I}

\page{25}
\note{Les six premières lignes sont annulées par une diagonale et cernées d'une accolade. La page reprend, en clair, la décomposition de $\mathcal{A}$ en facteurs $\mathcal{A}_i$ de type I, puis passe aux algèbres $\mathcal{A} \subset L(\mathcal{H})$ et à la formule $T(f) = \int_0^1 T(f(s))\,ds$.}

\struck{\ill{} $\mathcal{A}$ \uncertain{est} \uncertain{algèbre} \uncertain{de} \uncertain{Neumann} \ill{}, \uncertain{produit} \ill{} \uncertain{des} \uncertain{algèbres} $\mathcal{A}_i$ ($i \in I$) \uncertain{et} \uncertain{d'une} \uncertain{algèbre} \uncertain{de} $\mathcal{B}$ \ill{} \uncertain{les} \uncertain{traces} \ill{} $\mathcal{A}_i$ \uncertain{de} \uncertain{facteurs} \ill{}. \uncertain{Si} $T_i$, \uncertain{i.e.} $\lambda_i > 0$, \uncertain{il} \ill{} \uncertain{trace} \uncertain{usuelle}.} \uncertain{Mais} \uncertain{chaque} $\mathcal{A}_i$ \uncertain{permet} \ill{} \ill{} \uncertain{plongée} \uncertain{dans}, \uncertain{une} \uncertain{fois} \ill{} \uncertain{algèbre} \uncertain{de} \uncertain{Neumann} \ill{} \uncertain{l'algèbre} \ill{} $\mathcal{A}$ \ill{} \ill{} \uncertain{les} \uncertain{fonctions} \uncertain{caractéristiques} \ill{} \uncertain{dans} \uncertain{les} \uncertain{parties} \ill{} \uncertain{algèbre} \ill{} \uncertain{Neumann} \ill{} \uncertain{les} \ill{} \uncertain{de} $\mathcal{B}$, \uncertain{Donc} \ill{} \uncertain{que} \ill{} $T$ \ill{} \uncertain{commutant} \uncertain{des}, \uncertain{et} \uncertain{si} \ill{} \uncertain{une} \uncertain{trace} \uncertain{normale} \uncertain{fidèle} \ill{} \uncertain{On} \uncertain{peut} \uncertain{donc} \ill{} \uncertain{un} \uncertain{produit} \uncertain{de} $\mathcal{B}$ \ill{} \ill{} \uncertain{fini} \uncertain{fidèle}, \ill{} \uncertain{et} \uncertain{une} $\mathcal{A}$ \uncertain{de} \uncertain{type} \ill{} \uncertain{dénombrable}. \uncertain{L'algèbre} \uncertain{en} \uncertain{question}, \ill{} \uncertain{nouvelles} \uncertain{et} \uncertain{traces} \ill{} $\mathcal{A} \subset L(\mathcal{H})$, \uncertain{munie} \uncertain{des} (\uncertain{ou} $\mathcal{H} \otimes \mathcal{H}'$) \uncertain{avec} \uncertain{la} \uncertain{trace} \ill{} \uncertain{dans} \ill{} \uncertain{fonctions} \ill{} \uncertain{Neumann} \ill{} \uncertain{que} \ill{}

\uncertain{Il} \uncertain{est} \uncertain{aussi} \uncertain{à} \uncertain{remarquer} \uncertain{que} \ill{} \uncertain{fidèle} $T$, \uncertain{on} \uncertain{peut} \ill{} \uncertain{de} \uncertain{nouveau} \ill{} \uncertain{la} \uncertain{trace} \uncertain{normale}, \uncertain{à} \uncertain{savoir} \uncertain{celle} \ill{}. \uncertain{Par} \uncertain{exemple}, \uncertain{on} \uncertain{prend} \ill{} \struck{\uncertain{semi-finie}} \ill{} \uncertain{le} \uncertain{produit} \ill{} \uncertain{de} $[0,1]$ \uncertain{dans} $\mathcal{A}$ \ill{} \uncertain{et} \ill{} \uncertain{sur} $\mathcal{H}$ \ill{}, \uncertain{on} \uncertain{pose}
\[
  T(f) = \int_0^1 T(f(s))\, ds\,.
\]
\uncertain{Si} $\mathcal{B}$ \uncertain{est} \uncertain{une} \uncertain{algèbre} \ill{} \uncertain{munie} \uncertain{de} \uncertain{la} \uncertain{trace} $T$, \uncertain{et} $\mathcal{A}$ \uncertain{une} \uncertain{sous-algèbre} \ill{} \uncertain{de} \uncertain{Neumann}, \uncertain{alors} $\mathcal{A} \subset \mathcal{A}_2$ \struck{\ill{}} \ill{} \uncertain{dans} \uncertain{l'algèbre} \ill{} \uncertain{définie} \uncertain{de} \uncertain{la} \uncertain{fonction} \uncertain{de} \uncertain{la} \uncertain{trace} \uncertain{sur} $\mathcal{A}$, $\mu_A$ \uncertain{est} \uncertain{la} \uncertain{même} \uncertain{pour} \ill{} $\mathcal{A}$ \uncertain{en} \uncertain{réarrangement} \ill{} $\mathcal{A}$ \uncertain{ou} $\mathcal{B}$, \uncertain{d'où} \uncertain{les} \uncertain{questions} \ill{} \uncertain{les} \uncertain{propriétés} \ill{} \uncertain{on} \uncertain{voit} \ill{} \struck{\uncertain{les} \ill{}} \uncertain{peut} \ill{} \uncertain{plongé} \ill{} \uncertain{bien} \ill{} $\mathcal{A}$ \ill{} \ill{}
\note{le bas de la page reste illisible ; la phrase se termine au bord inférieur}

\page{26}
\emph{Prop 6} \struck{\uncertain{Soit} $A$ \ill{} \uncertain{hermitien} \ill{}} \note{cerné et biffé, au-dessus de la ligne} Soit $A \in \underline{b}^{+}$, \ill{}.
\[
  \text{(5)}\qquad \varphi_A^{-}(t) = \inf_{\substack{E \in \mathcal{P}(\mathcal{A})\\ \operatorname{Tr} E < t}}\ \sup_{x \in (1-E)\mathcal{H}} (Ax, x)
  = \inf_{\substack{E \in \mathcal{P}(\mathcal{A})\\ \operatorname{Tr} E < t}} \|(1-E)A(1-E)\|
\]
\note{$\mathcal{P}(\mathcal{A})$ est écrit d'un $P$ rond, pour l'ensemble des projecteurs de $\mathcal{A}$ ; « Tr » est écrit en toutes lettres ici, « $T(E)$ » dans les formules (6) et (7)}

\emph{Démonstration} Soit $\underline{C} = C_0(M)$ \note{cerné} \uncertain{une} \uncertain{sous-algèbre} \uncertain{commutative} \uncertain{maximale} \ill{} \uncertain{de} $\underline{b}$, \uncertain{contenant} $A$. \ill{} \uncertain{On} \uncertain{peut} \ill{} \uncertain{la} \uncertain{borne} \uncertain{dans} $\underline{C}$, \uncertain{i.e.} \uncertain{sur} $\underline{C}$, \uncertain{si} $E$ \uncertain{parcourt} \uncertain{les} \uncertain{projecteurs} \uncertain{de} $\underline{C}$, \uncertain{on} \uncertain{trouve}
\[
  \inf_{\substack{E \subset M\ \text{mesurable}\\ m(E) < t}} \|(1-E)A\|_\infty \geq \varphi_A^{-}(t)
\]
\note{le signe de comparaison est surchargé, un $=$ sous le $\geq$ ; le $A$ du premier membre est repassé} ($\uncertain{N°}~1$, \uncertain{formule} (4)) \uncertain{d'où} \uncertain{donc} \uncertain{dans} (5), $\varphi_A^{-}(t)$ \uncertain{est} $\geq$ \uncertain{le} \uncertain{second} \uncertain{membre} \ill{}.

\marginal{en marge gauche, sideways, à hauteur du corollaire 1 : « on pose $|A| = \sqrt{A^{*}A}$ »}

\emph{Corollaire 1} $A \in \underline{b}$ \uncertain{quelconque}, \ill{}
\[
  \text{(6)}\qquad \varphi_{|A|}^{-}(t) = \inf_{\substack{E \in \mathcal{P}(\mathcal{A})\\ T(E) < t}}\ \sup_{x \in (1-E)\mathcal{H}} \|Ax\|
  = \inf_{\substack{E \in \mathcal{P}(\mathcal{A})\\ T(E) < t}} \|A(1-E)\|
\]

\emph{Corollaire 2} $A \in \underline{b}$, $B \in \mathcal{A}$. \uncertain{Alors} \ill{}
\[
  \text{(7)}\qquad \left\{
  \begin{array}{l}
    \varphi_{|BA|} \leq \|B\|\, \varphi_{|A|} \\[0.5ex]
    \varphi_{|AB|} \leq \|B\|\, \varphi_{|A|}
  \end{array}\right.
\]
\note{le $B$ de la seconde ligne est surchargé}

\marginal{en marge gauche, sideways, à hauteur de (7) : $\varphi_{|A|}(\ldots) \leq \varphi(\ldots)$, \ill{} ; puis, plus bas : \uncertain{Cor.}~5 \ill{}}

\emph{Corollaire 3} $\varphi_A$ \uncertain{est} \uncertain{fonction} \uncertain{croissante} \uncertain{de} $A \in \underline{b}^{+}$.

\emph{Corollaire 4} \uncertain{Les} \uncertain{fonctions} $A \to \varphi_{|A|}^{-}(t)$ \uncertain{sur} $\underline{b}$ \uncertain{sont} \uncertain{semi-continues} \uncertain{supérieurement} \uncertain{quand} $\underline{b}$ \uncertain{est} \uncertain{muni} \uncertain{de} \uncertain{sa} \uncertain{topologie} \ill{}
\note{une accolade en marge gauche embrasse le corollaire ; les derniers mots sont perdus au bord de la page}

\page{27}
\emph{Prop 7} Soit $A \in \underline{b}$, alors
\[
  \mu_{|A|} = \mu_{|A^{*}|}, \qquad\text{i.e.}\qquad \varphi_{|A|} = \varphi_{|A^{*}|}
\]
\uncertain{On} \uncertain{a} $\varphi_{|A|}^{2} = \varphi_{A^{*}A}$, $\varphi_{|A^{*}|}^{2} = \varphi_{AA^{*}}$, \uncertain{il} \uncertain{suffit} \uncertain{de} \uncertain{prouver} $\varphi_{A^{*}A} = \varphi_{AA^{*}}$, \struck{i.e. $\mu_{AA^{*}} = \mu_{A^{*}A}$ \ill{}} \uncertain{prop.}~\uncertain{6}, \uncertain{un} \ill{} $4$, \ill{} \uncertain{les} \ill{} $2$ \uncertain{mais} \uncertain{dans} \struck{\uncertain{il} \uncertain{suffit} \ill{}} \uncertain{ceci} \uncertain{il} \uncertain{suffit} \uncertain{de} \uncertain{supposer} \struck{\ill{}} $A \in \underline{a}$, \uncertain{car} $\mu_{AA^{*}}$ \uncertain{et} \struck{$\mu_{AA^{*}}$} $\mu_{A^{*}A}$ \uncertain{sont} \ill{} \uncertain{sur} $\mathbf{R}^{+}$, \struck{\ill{}} \note{une ligne biffée entre crochets} \uncertain{et} \uncertain{il} \uncertain{suffit} \uncertain{de} \uncertain{prouver} $\mu_{AA^{*}}(x^{p}) = \mu_{A^{*}A}(x^{p})$ \uncertain{pour} $p > 0$, \ill{}, i.e.
\[
  T\bigl((AA^{*})^{p}\bigr) = T\bigl((A^{*}A)^{p}\bigr),
\]
\uncertain{ce} \uncertain{qui} \uncertain{résulte} \uncertain{aussitôt} \uncertain{de} $T(RS) = T(SR)$. \note{le second membre est repassé} \uncertain{Et} \uncertain{les} \ill{} \uncertain{peuvent} \uncertain{d'ailleurs} \uncertain{définir} \ill{} \uncertain{à} \uncertain{savoir} $\varphi_{|A|} = \varphi_{|A^{*}|}$ \uncertain{et} \uncertain{les} \ill{}, $\Phi_{|A|} = \Phi_{|A^{*}|}$, $\Psi_{|A|} = \Psi_{|A^{*}|}$, etc \struck{\ill{}} \ill{} \uncertain{l'indice} \uncertain{de} $A$.

\section{Compléments sur le déterminant}

\note{Titre de sa main, page 27 : « N° 6 » dans un cadre, puis « Compléments sur le déterminant » souligné. C'est le « N° 5 » du plan de la page 2, renuméroté 6 ici comme le plan lui-même l'avait d'abord fait.}

\struck{\uncertain{Définition}} \ill{}
\[
  \text{(1)}\qquad \Delta(A) = \struck{\ill{}}\ \exp\int_0^{T(1)} \log\varphi_{|A|} \qquad \text{\uncertain{pour} } A \in \tilde{\underline{b}} \text{ \uncertain{d'abord}, \uncertain{puis} } A \in \tilde{\underline{b}}
\]
\note{le second membre est écrit sous une première version biffée ; le « $\tilde{\underline{b}}$ » porte un tilde, et un mot devant lui est raturé ; ce que $\tilde{\underline{b}}$ désigne se déduit de la page 28 : les $\lambda 1 + B$, $B \in \underline{b}$}

\struck{\uncertain{Donc} \ill{}, $E$ \ill{} \uncertain{définition} \ill{}} \struck{$\Delta(A) = \Delta(|A|)$, \uncertain{et} \ill{} \uncertain{si} $A \in \underline{b}$ \ill{} $\Delta(|A|)$ \ill{}} \struck{\uncertain{et} \ill{}} $\Delta(A)$ \uncertain{est} \uncertain{défini} \ill{} \uncertain{pour} \ill{}, \uncertain{et} $0 \leq \Delta(A) \leq +\infty$. \uncertain{Donc}
\begin{itemize}
  \item[a)] \uncertain{Si} $T$ \uncertain{est} \uncertain{finie}, \ill{} \uncertain{pour} \uncertain{tout} $A \in \mathcal{A}$, \uncertain{et} $\Delta(\lambda A) = |\lambda|^{T(1)}\Delta(A)$, \struck{$\Delta(1) = 1$ \ill{}}
  \item[b)] \uncertain{Si} $T$ \uncertain{est} \uncertain{semi-finie} \ill{}, \uncertain{si} \ill{} $A \geq 0$, $A = \lambda 1 + B$ ($B \in \underline{b}$, $B$ \uncertain{hermitien} \ill{}). \uncertain{Si}
\end{itemize}
\note{la phrase se poursuit page 28}

\page{28}
\begin{itemize}
  \item[a.] \uncertain{Si} $|\lambda| < 1$, $\Delta(A) = 0$
  \item[b.] \uncertain{Si} $|\lambda| > 1$, $\Delta(A) = +\infty$
  \item[c.] \uncertain{Si} $|\lambda| = 1$, $0 \leq \Delta(A) \leq +\infty$ ; \struck{$\Delta(A) < \infty$} \uncertain{et} $\Delta(A) < +\infty \Longleftrightarrow B \in \underline{a}$
\end{itemize}
\note{les $|\lambda|$ sont surchargés ; on lit aussi $\|\lambda\|$}

\uncertain{Cela} \uncertain{résulte} \ill{} \uncertain{précédent}, \uncertain{à} \uncertain{savoir} \ill{} $0 \leq \Delta(1+A) < +\infty$ \uncertain{si} $A \in \underline{a}$, \uncertain{et} \ill{} \uncertain{particulier} \ill{} \uncertain{seulement} \uncertain{les} \uncertain{valeurs} \ill{} \uncertain{dans} $1 + \underline{a}$ \ill{} $\Delta(A)$, \uncertain{comme} \ill{} \uncertain{définie} \uncertain{à} \ill{} \uncertain{seulement} \struck{\ill{} $1 + B$, $B \in \underline{a}$}. \uncertain{Donc} \ill{} \uncertain{sur} $1 + \underline{a}$.

\emph{Prop 8}
\[
  \left\{
  \begin{array}{l}
    \Delta(A) = \Delta(A^{*}) \\
    \Delta(AB) = \Delta(A)\,\Delta(B) \\
    \Delta(1) = 1
  \end{array}\right.
  \qquad A, B \in 1 + \underline{a}
\]
\note{un long trait courbe enferme l'énoncé et le relie à « $A, B \in 1 + \underline{a}$ » ; à droite de la première ligne, un mot raturé}

\struck{Si $A \in \mathcal{A}$, \ill{} \uncertain{un} $A\mathcal{H}$, \ill{}. $A \in \underline{a}^{+}$, \uncertain{alors}}
\[
  \text{(5)}\qquad \struck{\Delta(1 + A) = \Delta_E(1_E + A_E)} \qquad \struck{(E = \text{\uncertain{support}} \ill{}}
\]
\struck{\uncertain{désigne} $E$ \uncertain{et} $A$ \uncertain{un} \uncertain{opérateur} \ill{} \uncertain{comme} \ill{} $A_E$ \uncertain{et} $1_E$ \uncertain{désignent} \ill{} \uncertain{de} \uncertain{l'algèbre} \uncertain{de} \uncertain{von} \uncertain{Neumann} $\mathcal{A}_E = E\mathcal{A}E$, \uncertain{et} $\Delta_E$ \uncertain{le} \uncertain{déterminant} \ill{} \uncertain{dans} $\mathcal{A}_E$ \uncertain{muni} \uncertain{de} \uncertain{la} \uncertain{trace} \uncertain{induite} \ill{} \uncertain{de} $\mathcal{A}$.}
\note{le bloc, de « Si $A \in \mathcal{A}$ » à « de $\mathcal{A}$ », est annulé par deux longues diagonales ; il est repris page 29 sous le numéro (5'), avec la notation $\Delta^{E}$ en exposant qui s'installe dès les lignes suivantes}

\uncertain{Soit} $E \in \mathcal{P}(\mathcal{A})$, \uncertain{posons} $\mathcal{A}_E = E\mathcal{A}E$ (\uncertain{sous-algèbre} \ill{} \uncertain{de} $\mathcal{A}$, \uncertain{ayant} \uncertain{l'unité} $E$, \uncertain{et} \ill{} \uncertain{une} \uncertain{algèbre} \uncertain{de} \uncertain{von} \uncertain{Neumann} \uncertain{dans} $E\mathcal{H}$, \struck{\ill{}} \uncertain{les} \ill{} \uncertain{la} \uncertain{restriction} \uncertain{de} \uncertain{la} \uncertain{trace} $T$ \note{en interligne : « $T_E$ »} \uncertain{est} \uncertain{une} \uncertain{trace} \uncertain{normale} \uncertain{semi-finie} \uncertain{fidèle}, \uncertain{et} \uncertain{soit} $\Delta^{E}(A)$ \uncertain{pour} $A \in \mathcal{A}_E$ \uncertain{le} \uncertain{déterminant} \uncertain{de} $A$ \uncertain{considéré} \uncertain{dans} $\mathcal{A}_E$ \uncertain{relativement} \ill{}

\page{29}
\uncertain{comme} \uncertain{un} \uncertain{déterminant} \ill{} \uncertain{relatif} \uncertain{à} $\mathcal{A}_E$. \uncertain{On} \uncertain{a} \uncertain{alors}
\[
  \text{(5')}\qquad \Delta(1 + A) = \Delta^{E}(1_E + A) \qquad \text{si } \struck{A \in \mathcal{A}_E}\ A \in \mathcal{A}_E
\]
\note{le numéro est un « (5) » surchargé d'un prime ; l'exposant $E$ de $\Delta$ est écrit sur un indice raturé, et le $1_E$ sur un $A$ ; le « $A \in \mathcal{A}_E$ » est réécrit sous une première version biffée}
\struck{\uncertain{et} \uncertain{si} \ill{} (\uncertain{quand} \ill{}) \ill{} \uncertain{de} \ill{} \uncertain{parcourt} \ill{} $1 + $ \uncertain{opérateur} \ill{} \uncertain{tel} \ill{} \uncertain{fini}} \note{trois lignes biffées et cernées} \uncertain{comme} \uncertain{on} \uncertain{le} \uncertain{vérifie} \uncertain{aussitôt}.

\emph{Prop 8}
\[
  \left\{
  \begin{array}{l}
    \Delta(A) = \Delta(A^{*}) \\
    \Delta(AB) = \Delta(A)\,\Delta(B) \\
    \Delta(1) = 1
  \end{array}\right.
  \qquad \text{si } A, B \in 1 + \underline{a}
\]
\note{une flèche relie l'accolade à la seconde ligne}

\uncertain{Seule} \uncertain{la} \uncertain{seconde} \uncertain{formule} \ill{} \uncertain{démonstration}. \uncertain{Par} \uncertain{raison} \uncertain{de} \uncertain{continuité} \ill{} \uncertain{on} \uncertain{peut} \uncertain{supposer} \uncertain{qu'il} \uncertain{existe} $E \in \underline{a}$ \ill{} \uncertain{un} \ill{} \uncertain{projecteur} \ill{}, \uncertain{tel} \uncertain{que} \ill{} $A \in \mathcal{A}_E$, $B \in \mathcal{A}_E$, \uncertain{et} \ill{} \uncertain{dans} $\mathcal{A}_E$ \struck{$\mathcal{A}$ \ill{} \uncertain{à} \ill{}} \uncertain{on} \uncertain{peut} \uncertain{supposer} \ill{} \uncertain{le} \uncertain{cas} \ill{} \uncertain{finie}. \struck{\ill{} \uncertain{les} \ill{}} \uncertain{de} \uncertain{plus} (\uncertain{par} \uncertain{raison} \uncertain{de} \uncertain{continuité}) \ill{}

\struck{$A$ \uncertain{et} $B$ \ill{} \uncertain{peut} \ill{}}
\begin{align*}
  \struck{\Delta(A)^{2}} &\struck{= \Delta(AA^{*})} \\
  \struck{\Delta(A^{*})^{2}} &\struck{= \Delta(A^{*}A)} \\
  \struck{\Delta(AB)^{2}} &\struck{= \Delta(B^{*}A^{*}AB) = \Delta(B^{*}HB) \overset{?}{=} \Delta(\uncertain{B^{*}A})\,\Delta(B^{*}B)}
\end{align*}
\struck{\uncertain{Soit} $B^{*}B = |B|^{2}$ ; \uncertain{on} \uncertain{a} $\Delta(B^{*}HB) = \Delta(LHL)$ \uncertain{On} \uncertain{est} \uncertain{ramené} \ill{} \uncertain{pour} $\Delta(LHL)$}
\note{tout ce bloc, de « $A$ et $B$ » à « $\Delta(LHL)$ », est enfermé dans de grandes boucles qui l'annulent ; le point d'interrogation sur le dernier signe $=$ est de sa main. $H$ est mis pour $A^{*}A$, $L$ pour $|B|$ ; la démonstration est reprise page 30 par la décomposition polaire. Le reste de la page est blanc.}

\page{30}
\struck{Soit} \uncertain{Nous} \uncertain{avons} \uncertain{à} \uncertain{définir} \ill{}. \uncertain{Supposons} \uncertain{que} $\mathcal{A}$ \uncertain{soit} \uncertain{à} \uncertain{trace} \uncertain{finie}. \uncertain{Soient} $X, Y$ \uncertain{deux} \uncertain{projecteurs} \uncertain{équivalents} \uncertain{dans} $\mathcal{A}$, \uncertain{i.e.} \uncertain{il} \uncertain{existe} \ill{} \uncertain{transformé} \ill{} $U \in \mathcal{A}$ \uncertain{avec}
\[
  U^{*}U = X \qquad UU^{*} = Y
\]
Soit $A \in \mathcal{A}$ \uncertain{tel} \uncertain{que} $AX(\mathcal{H}) \subset Y(\mathcal{H})$, i.e. $\sigma(AX) \leq Y$. \uncertain{Considérons} \struck{\ill{}} \uncertain{posons} \ill{}. \uncertain{Alors}
\[
  \Delta^{X}(U^{*}AX) = \Delta^{Y}(YAU^{*})
\]
\note{ici et dans (5) et (6), l'exposant est écrit sur un indice raturé : la notation $\Delta_X$ est corrigée en $\Delta^{X}$} \uncertain{Ce} \uncertain{qui} \uncertain{se} \uncertain{voit} \uncertain{facilement} \ill{} \uncertain{car} \ill{} \uncertain{ne} \uncertain{dépend} \uncertain{que} \uncertain{de} \ill{} \uncertain{et} \uncertain{on} \uncertain{peut} \uncertain{donc} \uncertain{remplacer} $U$ \uncertain{par} \uncertain{un} $V$ \uncertain{tel} \uncertain{que} \ill{}, \uncertain{tel} \uncertain{que} \ill{}, \uncertain{est} \uncertain{unitairement} \ill{} $U$ \ill{} $V^{*}V = X$, $VV^{*} = Y$. \uncertain{On} \uncertain{pose} \uncertain{donc}
\[
  \text{(5)}\qquad \Delta^{X,Y}(A) = \Delta^{X}(U^{*}AX) = \Delta^{Y}(YAU^{*}) \qquad (\sigma(AX) \leq Y)
\]
\uncertain{Soit} \uncertain{de} \uncertain{plus} $Z$ \uncertain{un} \uncertain{projecteur} \uncertain{équivalent} \uncertain{à} \ill{} $Y$, \uncertain{et} $B \in \mathcal{A}$ \uncertain{tel} \uncertain{que} $\sigma(BY) \leq Z$. \uncertain{On} \uncertain{a} \uncertain{alors} \uncertain{les} \uncertain{formules} \ill{} \uncertain{prop.}~8 \ill{}
\[
  \text{(6)}\qquad \Delta^{X,Z}(BA) = \Delta^{X,Y}(A)\,\Delta^{Y,Z}(B)
\]
\uncertain{et} \uncertain{ces} \uncertain{deux} \uncertain{propriétés} \uncertain{ne} \ill{} $A$ \ill{} $=$ \ill{}

\emph{Prop 9}
\[
  \Delta_A(t) = \sup_{\operatorname{Tr} E \leq t} \Delta^{E}(EAE) \qquad (A \in \underline{b})\ (t \leq \operatorname{Tr} 1)
\]
\note{l'énoncé est cerné d'un trait ; le « Tr » est en toutes lettres, et le second membre porte encore l'indice raturé sous l'exposant $E$} \uncertain{Comme} \uncertain{les} \uncertain{deux} \uncertain{membres} \ill{} \uncertain{continus} \ill{} \uncertain{en} $A$ \ill{} \uncertain{restreint} \ill{} \uncertain{modulo} \ill{} \uncertain{commutative} \ill{} $\underline{C} \subset \underline{b}$, $\underline{C} = C_0(M)$ ; $|M| = T(1)$. (\uncertain{diagonalisation} \ill{} \uncertain{par} \ill{}) \ill{} \uncertain{on} \uncertain{a} \ill{}
\[
  \Delta_A(t) \geq \sup_{\operatorname{Tr} E \leq t} \Delta_E(EAE)
\]
\note{ici l'indice $E$ n'est pas corrigé en exposant} \ill{} \uncertain{réciproquement}, \uncertain{il} \uncertain{faut} \uncertain{prouver} : $\Delta^{E}(EAE) \leq \Delta_A(t)$ \uncertain{si} $T(E) \leq t$. \uncertain{Or} \ill{}. \struck{$\Delta^{E}(EAE) = \Delta^{E}_{EAE}(t)$}
\note{la démonstration se poursuit page 31}

\page{31}
\[
  \Delta^{E}(EAE) = \Delta^{E}_{EAE}(T(E)) = \Delta_{EAE}(t)\struck{\ill{}} \leq \Delta_A(t)\struck{\ill{}}
\]
\note{deux termes raturés, après $\Delta_{EAE}(t)$ et en fin de ligne} (\uncertain{car} \ill{} \uncertain{on} \uncertain{a} $\varphi_{|EAE|} \leq \|E\|\,\|E\|\,\varphi_{|A|} = \varphi_{|A|}$) \note{la parenthèse est cernée}

\emph{Corollaire} \uncertain{Si} $E \simeq F$, \struck{\ill{}} \uncertain{et} $A \in \mathcal{A}$, \uncertain{on} \uncertain{a}
\[
  \Delta^{EF}(FAE) \leq \Delta_A(t) \qquad (t = T(E))
\]
\struck{\emph{Corollaire} \ill{} $\Delta^{E}$ \ill{}} \struck{(8) $\Delta_{AB}(t) \leq \Delta_A(t)\,\Delta_B(t)$ \ill{} $t$.} \note{un premier énoncé de (8), cerné, est biffé ici et repris comme théorème}

\emph{Théorème 1} \uncertain{On} \uncertain{a}
\[
  \Delta_{AB}(t) \leq \Delta_A(t)\,\Delta_B(t) \qquad (A, B \in \underline{b},\ \struck{t \leq |M|}\ t \geq 0)
\]
\uncertain{Comme} \uncertain{les} \uncertain{deux} \uncertain{membres} \uncertain{sont} \uncertain{fonctions} \uncertain{continues} \ill{} \uncertain{de} $A$ \ill{} \uncertain{à} \ill{} \uncertain{techniques}, \uncertain{et} $T(1) < +\infty$. \uncertain{Comme} \note{cerné}
\[
  \Delta_{AB}(t) = \sup_{T(E) \leq t} \Delta_E(EABE)
\]
\uncertain{il} \uncertain{suffit} \uncertain{de} \uncertain{prouver} \uncertain{que} \uncertain{pour} $T(E) = t$ \uncertain{on} \uncertain{a}
\[
  \Delta^{E}(EABE) \leq \Delta_A(t)\,\Delta_B(t)\,.
\]
\struck{\uncertain{Posons} \ill{} \ill{} $\sigma(BE)$ \ill{}} \note{deux lignes biffées, la seconde cernée} \uncertain{Soit} $F$ \uncertain{un} \uncertain{projecteur} \ill{} $\geq \sigma(BE)$, \uncertain{et} $\simeq E$, \uncertain{soit} $A' = EAF$, $B' = FBE$ \struck{$= $ \ill{}}, \ill{}. \struck{$\sigma(A'E) = \sigma(BE) \leq F$, $\sigma(B'F) = $} $A'B' = EABE$, \uncertain{et} \ill{} \uncertain{en} \uncertain{vertu} \uncertain{de} (6),
\[
  \Delta^{E}(EABE) = \Delta^{E}(A'B') = \Delta^{E,F}(B')\,\Delta^{F,E}(A')\,.
\]
\note{l'exposant du second $\Delta$ est écrit sur un indice raturé ; les deux facteurs du dernier membre portent les exposants « $E,F$ » et « $FE$ »}
\struck{\uncertain{Or} $\bigl(\Delta^{EF}(B')\bigr)^{2} = \Delta^{E}(EB^{*}FBE) = \Delta^{E}_{EB^{*}FBE}(t)$ ; $\bigl(\Delta^{FE}(A')\bigr)^{2} = \Delta^{E}(EA^{*}FAE) = \ill{}$ ; \uncertain{prop.}~9) \uncertain{Or} $\Delta_{B^{*}FB}(t) = \Delta_{B^{*}F}(t) \leq \Delta_{B^{*}}(t)$ ; \uncertain{première} \uncertain{égalité} \uncertain{résulte} \uncertain{d'une} \uncertain{définition}, \uncertain{la} \uncertain{deuxième} \uncertain{résulte} \uncertain{de} \uncertain{prop.}~6, \uncertain{cor.}~2)}
\note{le bloc est enfermé dans trois boucles emboîtées qui l'annulent ; on y lit encore « $\Delta_{B^{*}FB}(t) = \Delta_{B^{*}F}(t) \leq \Delta_{B^{*}}(t)$ » et le renvoi à « prop 6, cor. 2 »}
\uncertain{Or} $\Delta_{B'}(t) = \Delta_B(t)$. \ill{} \uncertain{Comme} \ill{} \uncertain{on} \uncertain{a} \ill{}
\begin{align*}
  \Delta^{E,F}(\uncertain{FBE}) &\leq \Delta_B(t) \\
  \Delta^{F,E}(EAF) &= \Delta_A(t)
\end{align*}
\note{l'argument du premier $\Delta$ est surchargé ; le signe de la seconde ligne, $=$, est ce que la page porte} \ill{} \uncertain{résulte} \uncertain{du} \uncertain{corollaire} \uncertain{précédent}, \uncertain{quand} \uncertain{on} \uncertain{lit} \ill{} \uncertain{prop.}~9 :
\note{la démonstration se poursuit page 32}

\page{32}
\marginal{en marge gauche, sideways : « corollaire prop 9 »}

Soit $E \simeq F$ ($E, F$ \uncertain{projecteurs} \uncertain{équivalents}), $T(E) = t$, \uncertain{et} $A \in \mathcal{A}$, \uncertain{alors}
\[
  \Delta^{EF}(FAE) \leq \Delta_A(t)
\]
\uncertain{En} \uncertain{effet}, \uncertain{on} \uncertain{a}
\[
  \Delta^{EF}(FAE) = \Delta^{E}(U^{*}FAE) = \Delta^{E}(U^{*}AE) \leq \Delta_{U^{*}A}(t) \leq \Delta_A(t) \quad \text{\uncertain{car}}\ \|U^{*}\| \leq 1
\]
(\uncertain{on} \uncertain{utilise} \ill{} \uncertain{prop.}~6, \uncertain{cor.}~2)
\note{une accolade en marge gauche embrasse l'énoncé et sa démonstration}

\section{Inégalités fondamentales}

\note{Titre de sa main, souligné, page 32, sans numéro. C'est le « N° 6 » du plan de la page 2, « Les inégalités pour $AB$ », que la renumérotation du N° 5 en N° 6 a décalé.}

\uncertain{Le} \uncertain{Th.}~1 \uncertain{peut} \uncertain{s'écrire} \uncertain{encore} :
\[
  \text{(1)}\qquad \Delta_{\varphi_{|AB|}} \leq \Delta_{\varphi_{|A|}\varphi_{|B|}}
\]
(\uncertain{car} $\Delta_{fg} = \Delta_f\,\Delta_g$ \uncertain{pour} \uncertain{des} \uncertain{fonctions} \uncertain{décroissantes} \ill{} \ill{} \uncertain{sur} $\mathbf{R}^{+}_{*}$) \uncertain{Par} \uncertain{suite}, \uncertain{utilisant} \uncertain{N°}~\uncertain{4}, \uncertain{corollaire}~1, \ill{}

\emph{Th. 2} \uncertain{Si} $A, B \in \underline{b}$, \uncertain{et} \uncertain{si} $W$ \uncertain{est} \uncertain{une} \uncertain{fonction} \uncertain{de} Weyl \uncertain{croissante} \ill{} \uncertain{dans} $K^{\infty}(]0, T(1)])$, \ill{}
\[
  \text{(2)}\qquad W(AB) \leq W(\varphi_{|A|}\varphi_{|B|})
\]
\note{une accolade en marge gauche embrasse l'énoncé ; l'intervalle $]0, T(1)]$ est surchargé}

\emph{Corollaire 1} \uncertain{On} \uncertain{a} \struck{$\Phi_{|AB|}(t) \leq$}
\[
  \text{(3)}\qquad \int_0^{t} \varphi_{|AB|}(s)\, ds \leq \int_0^{t} \varphi_{|A|}(s)\,\varphi_{|B|}(s)\, ds
\]
\uncertain{pour} \struck{\ill{}} $t \geq 0$. \uncertain{En} \uncertain{particulier}
\[
  \text{(4)}\qquad \boxed{\ \int \varphi_{|AB|} \leq \int \varphi_{|A|}\,\varphi_{|B|}\ }
\]
\uncertain{d'où}
\[
  \text{(5)}\qquad |\operatorname{Tr} AB| \leq \struck{\ill{}}\ \operatorname{Tr}|AB| = \int \varphi_{|AB|} \leq \int \varphi_{|A|}\,\varphi_{|B|}
\]
\note{une accolade en marge gauche embrasse le corollaire de (3) à (5) ; (4) est encadrée d'un double trait}
(\uncertain{La} \uncertain{plupart} \uncertain{des} \uncertain{inégalités} \uncertain{importantes} \uncertain{peuvent} \uncertain{déjà} \uncertain{être} \uncertain{dites} \uncertain{conséquences} \uncertain{immédiates} \uncertain{de} (4).)

\page{33}
\uncertain{Soient} $p, q, r > 0$ \uncertain{avec} $\frac{1}{p} + \frac{1}{q} = \frac{1}{r}$, \uncertain{prenons} \ill{} $W(f) = \int |f|^{r}$, \uncertain{d'où} \uncertain{alors} \struck{$W(A)$} $W(A) = \operatorname{Tr}(|A|^{r})$, \struck{$W(AB)$ \ill{}} \uncertain{d'où} \ill{} \uncertain{alors} :
\[
  \operatorname{Tr}|AB|^{r} \leq \int \bigl(\varphi_{|A|}\varphi_{|B|}\bigr)^{r} \qquad\text{i.e.}\qquad \|AB\|_r \leq \|\varphi_{|A|}\varphi_{|B|}\|_r
\]
\struck{\ill{}} \uncertain{d'où} \ill{} \uncertain{en} \uncertain{vertu} \uncertain{de} \uncertain{l'inégalité} \uncertain{de} Hölder \uncertain{classique} :

\emph{Corollaire 2}
\[
  \text{(6)}\qquad \|AB\|_r \leq \|A\|_p\,\|B\|_q \qquad \Bigl(\frac{1}{r} = \frac{1}{p} + \frac{1}{q}\Bigr)
\]
(\struck{\ill{}} \uncertain{Inégalité} \uncertain{de} Hölder \uncertain{pour} \uncertain{opérateurs})

\emph{Corollaire 3} \uncertain{Soit} $A \in \mathcal{A}$ \ill{}, \ill{} \uncertain{de} $\mathcal{A}$ \uncertain{tel} \uncertain{que} \ill{}. \uncertain{Pour} $A \in \underline{b}$ \struck{\ill{}}
\[
  \text{(7)}\qquad \Phi_{|A|}(t) = \sup_{\sigma(B) \leq t} |\operatorname{Tr} AB| = \sup_{\substack{\operatorname{Tr}(\sigma(B)) \leq t\\ \|B\| \leq 1}} |\operatorname{Tr} BA|
\]
\note{le premier sup est barré d'un trait oblique qui le relie au second ; la condition « $\sigma(B) \leq t$ » est écrite sous « $\operatorname{Tr}$ » raturé} \uncertain{la} \uncertain{borne} \uncertain{sup} \uncertain{étant} \uncertain{prise} \uncertain{sur} \uncertain{les} $B$ \uncertain{tels} \uncertain{que} \ill{} \uncertain{ces} \uncertain{conditions}. \ill{} \uncertain{Le} \ill{} \uncertain{résulte} \uncertain{aussitôt} \uncertain{de} \uncertain{ce} \uncertain{qui} \uncertain{précède} \ill{} \uncertain{la} \uncertain{dernière} \uncertain{inégalité} ; \uncertain{d'autre} \uncertain{part}, \uncertain{si} $A = U|A|$, $|A| = U^{*}A$, \uncertain{on} \ill{} \uncertain{diagonalise} \ill{} \uncertain{trouve} $E$ \uncertain{projecteur} \uncertain{commutant} \uncertain{avec} \ill{}, \uncertain{avec} $\operatorname{Tr} E = t$, \uncertain{tel} \uncertain{que}
\[
  \operatorname{Tr} E|A| = \Phi_{|A|}(t), \quad \text{\uncertain{d'où}}\quad |\operatorname{Tr} EU^{*}A| = \Phi_{|A|}(t), \quad \text{\uncertain{et}}\quad \operatorname{Tr}\bigl(\sigma(EU^{*})\bigr) \leq \operatorname{Tr} E \leq t,\quad \|EU^{*}\| \leq 1.
\]
\uncertain{D'où} \ill{} (7), \uncertain{l'inégalité} \ill{} $\leq$, \uncertain{il} \ill{} \uncertain{prouver} \ill{}, \uncertain{on} \uncertain{a} \struck{$\varphi_{|B|} \leq$} $\varphi_{|B|} = \varphi_{|EB|}$ (\uncertain{si} $E = \sigma(B)$) \ill{} $\leq \|B\|\,\varphi_{|E|} = \varphi_{|E|}$, \uncertain{d'où} \uncertain{par} \ill{} ((5)) :
\[
  |\operatorname{Tr} BA| \leq \int \varphi_E\,\varphi_{|A|} = \int_0^{t} \varphi_{|A|}(s)\, ds = \Phi_{|A|}(t)\,.
\]
\uncertain{d'où} (7), \uncertain{cqfd}.

\emph{Corollaire 4} \uncertain{Les} \uncertain{fonctions} $A \to \Phi_{|A|}(t)$ \ill{} \uncertain{sont} \ill{} \uncertain{sur} $\underline{b}$, \ill{} \struck{$\Phi_{|A+B|}(t) =$} \ill{} \uncertain{sous-additives} :
\[
  \text{(8)}\qquad \Phi_{|A+B|} \leq \Phi_{|A|} + \Phi_{|B|}
\]
\note{la page s'arrête sur (8) ; la page 34 reprend le N° 4 depuis sa proposition 1, dans une seconde mouture}

\section{Réarrangements spectraux d'opérateurs, seconde mouture}

\note{Titre de l'éditeur. La page 34 reprend, sur un feuillet plus petit et d'une main plus serrée, les propositions 6 et 7 du N° 4 (pages 26-27) sous les numéros 1 et 2, avec leurs corollaires, puis (page 35) la proposition 9 du N° 6 ; la page 34 s'ouvre sur le mot « Astuce », souligné. On transcrit cette mouture à part, sans la fondre dans la première.}

\page{34}
\struck{\uncertain{On} \uncertain{suppose} \ill{} \uncertain{trace} \ill{} \uncertain{continue}.}

\emph{Astuce} : \uncertain{on} \uncertain{peut} \ill{} \ill{} \uncertain{une} \ill{}, \uncertain{et} \ill{} \uncertain{trace} \ill{}, \ill{} \uncertain{de} \ill{} \uncertain{une} \uncertain{trace} \uncertain{continue}.

\struck{\uncertain{Soit} \ill{} $\varphi_A(t)$ \ill{}}

\emph{Prop 1}
\[
  \varphi_A(t+0) = \inf_{\substack{E \in \mathcal{P}\\ \operatorname{Tr} E < t}}\ \sup_{x \in (1-E)\mathcal{H}} (Ax, x) \qquad (\text{\uncertain{trace} \uncertain{continue}})
\]
\note{le « $t+0$ » remplace le $\varphi_A^{-}$ de la page 26}

\uncertain{On} \uncertain{a} \uncertain{en} \uncertain{fait}, \uncertain{comme} \ill{} (\uncertain{partie}~2) $\varphi_A(t) \geq$ \uncertain{second} \uncertain{membre} \uncertain{parce} \ill{} \uncertain{on} \uncertain{voit} \uncertain{aussitôt} (\uncertain{diagonalisation} \ill{}) \ill{}, \uncertain{il} \uncertain{existe} \uncertain{donc} \uncertain{un} $E$ \uncertain{avec} $\operatorname{Tr} E < t$ \ill{}
\[
  \sup_{x \in (1-E)\mathcal{H}} (Ax, x) \geq \varphi_A^{-}(t)
\]
\note{cerné} \ill{} \uncertain{on} \uncertain{se} \uncertain{ramène} \uncertain{au} \uncertain{cas} \uncertain{où} $A$ \uncertain{est} \uncertain{régulier} : $\lambda F$, $F$ \uncertain{projecteur} \ill{} (\uncertain{i.e.} $A \geq \varphi_A^{\struck{-}}(t)\, F$, \uncertain{avec} $\operatorname{Tr} F \geq t$) \note{l'exposant de $\varphi_A$ est noirci} \ill{} \uncertain{Il} \uncertain{existe} \uncertain{donc} \ill{} \ill{}, \uncertain{si} $\operatorname{Tr} E < \operatorname{Tr} F$, \uncertain{alors} $(1-E)\mathcal{H} \wedge F\mathcal{H} \neq (0)$. \ill{} : \uncertain{Soit} $\mathcal{X} = E\mathcal{H}$, $\mathcal{Y} = F\mathcal{H}$, $\mathcal{X}' = \mathcal{H} \ominus \mathcal{X}$, $\mathcal{Y}' = \mathcal{H} \ominus \mathcal{Y}$ : \ill{} \uncertain{tel} \uncertain{que} \ill{}
\[
  \mathcal{X} \ominus (\mathcal{X} \wedge \mathcal{Y}') \sim \mathcal{Y} \ominus (\mathcal{Y} \wedge \mathcal{X}') \qquad \text{\uncertain{donc}}\quad \mathcal{Y} \wedge \mathcal{X}' = 0
\]
\uncertain{indiquerait} \struck{$\mathcal{X} \ominus$} $\mathcal{Y} \uncertain{\prec} \mathcal{X}$ \uncertain{d'où} $\operatorname{Tr} F \leq \operatorname{Tr} E$, \uncertain{absurde}.

\marginal{en marge gauche, sideways, sur trois lignes : \uncertain{si} $\varphi(t)$ \ill{}, $\varphi_{|1+A|} \leq 1 + \varphi_{|A|}$ \ill{} ; $\varphi_{1+\uncertain{|A|}} = 1 + \varphi_{|A|}$, \uncertain{si} \ill{} ; plus bas, souligné : « \emph{Corollaire 4} » puis « $\varphi_{\alpha(A)} = \alpha\,\varphi_A$ », un gribouillis raturé au-dessus}

\emph{Corollaire 1} $A$ \uncertain{quelconque} \ill{} :
\[
  \varphi_{|A|}(t+0) = \inf_{\substack{E \in \mathcal{P}\\ \operatorname{Tr} E < t}}\ \sup_{x \in (1-E)\mathcal{H}} \|Ax\| \qquad (\text{\uncertain{trace} \uncertain{continue}})
\]

\emph{Corollaire 2}
\[
  \varphi_{|PA|}\struck{(t)} \leq \|P\|\,\varphi_{|A|}\struck{(t)} \qquad (\text{\uncertain{trace} \uncertain{quelconque} \ill{}}), \qquad
  \varphi_{|AQ|} \leq \|Q\|\,\varphi_{|A|}
\]
\struck{\ill{} \uncertain{définit} \ill{}} \note{une ligne biffée, à droite de la seconde inégalité}

\emph{Corollaire 3} $\varphi_A$ \uncertain{croît} \uncertain{avec} \uncertain{l'hermitien} \ill{} \uncertain{pour} $A$ \ill{} \struck{\ill{}}

\emph{Corollaire 5} \uncertain{semi-continuité} \ill{} \uncertain{des} \uncertain{fonctions} $A \to \varphi_{|A|}(t)$ \ill{} \struck{$\varphi_{|A|}(t+0)$} \ill{} \uncertain{fonction} \uncertain{de} $A$.
\note{le corollaire 4 est celui de la marge ; une accolade relie les corollaires 3 et 5}

\emph{Prop 2} $\varphi_{|A|} = \varphi_{|A^{*}|}$ \uncertain{Il} \uncertain{suffit} \uncertain{de} \uncertain{prouver} \uncertain{que} $\varphi_{AA^{*}} = \varphi_{A^{*}A}$ \ill{} \uncertain{Il} \uncertain{suffit} \uncertain{de} \uncertain{prouver} \uncertain{pour} $A \in \underline{a}$ ! \uncertain{Pour} \ill{} \struck{$A \in \underline{a}$, \ill{}} \ill{} \uncertain{que} \ill{} \struck{$\mu_{|A|} = \mu_{|A^{*}|}$} \ill{}
\note{la phrase se poursuit page 35}

\page{35}
\[
  \mu_{|AA^{*}|}(x^{p}) = \mu_{|A^{*}A|}(x^{p}) \qquad \text{\uncertain{pour} \uncertain{tout} \ill{} \uncertain{entier} } p \geq 1,\ \text{\uncertain{soit}}
\]
\[
  \operatorname{Tr}\,\underbrace{AA^{*} \cdots AA^{*}}_{p} = \operatorname{Tr}\,\underbrace{A^{*}A \cdots A^{*}A}_{p}
\]
: \uncertain{il} \uncertain{suffit} \uncertain{de} \uncertain{prouver} \uncertain{pour} $A$ \ill{} \ill{} : 2 \ill{}.

\struck{\uncertain{Proposition}} \uncertain{Posons}
\[
  \Psi_{|A|}(t) = \Delta_t(A)
\]
\note{« Posons » est écrit sur « Proposition » biffé ; la notation $\Delta_t(A)$, indice $t$, remplace ici le $\Delta_A(t)$ des pages 22-33}

\emph{Proposition} (\uncertain{si} \uncertain{trace} \uncertain{continue})
\[
  \left\{
  \begin{array}{l}
    \Delta_t(A) = \displaystyle\sup_{\operatorname{Tr} E < t} \struck{\det}\ \Delta_E\, EAE \\[2ex]
    \Delta_t(1 + |A|) = \displaystyle\sup_{\substack{\operatorname{Tr}|B| < t\\ \|B\| \leq 1}} \struck{\det}\ \Delta_E\bigl(1_E + BA\uncertain{B}\bigr)
  \end{array}\right.
\]
\note{les deux « det » sont biffés et remplacés par $\Delta_E$ ; le dernier facteur de la seconde ligne est noirci. À droite : « \struck{$A$ \uncertain{quelconque}} (\uncertain{ou} \uncertain{bien} \uncertain{sup} \uncertain{borné} : $B \in \mathcal{P}$ \uncertain{si} $A \geq 0$) $\leq$ » ; un gribouillis précède l'accolade}

\uncertain{On} \uncertain{a} (\uncertain{décomposition} \uncertain{spectrale}) \uncertain{ce} \uncertain{signe} \ill{} :
\begin{align*}
  \Delta_E\, EAE &\leq \Delta_t(A) \qquad \text{\uncertain{si} } \operatorname{Tr} E \leq t \\
  \Delta_E(1_E + BAB) &\leq \Delta_t(1 + |A|)
\end{align*}
\note{les deux $\Delta_E$ sont écrits sur des « det » biffés ; un trait vertical sépare « si $\operatorname{Tr} E \leq t$ »} \uncertain{Il} \uncertain{faut} \uncertain{donc} \uncertain{prouver} \ill{}
\begin{align*}
  \Delta_E A' &\leq \Delta_t A' \\
  \Delta_E(1 + B') &\leq \Delta_t(1 + |B'|)
\end{align*}
\uncertain{soit} $A' = EAE$, \ill{}, $B' = BA$ \ill{} $\Delta_t(B') \leq \Delta_t(A)$ \uncertain{et} $\Delta_t(1 + |B'|) \leq \Delta_t(1 + |A|)$ \struck{$\leq \Delta_t(1+A)$} \uncertain{grâce} \uncertain{à} : \uncertain{Prop.}~1, \uncertain{corollaire}~2 (\uncertain{d'où} \uncertain{la} \uncertain{deuxième} \ill{}, \ill{} \uncertain{remarquer} \ill{} \uncertain{qui} \ill{} \uncertain{on} \uncertain{a} \ill{} \uncertain{et} \uncertain{finies}, \uncertain{donc} \ill{}, \uncertain{si} $\varphi_{1+A}$ \uncertain{existe}, \uncertain{alors}
\[
  \varphi_{E(1+A)E} \leq \varphi_{1+A}\ )
\]
\uncertain{Car} \struck{\ill{} \uncertain{réunion} \ill{}} \struck{$\det$} \struck{$\Delta B$} $\leq \Delta_t B$ \uncertain{si} $\operatorname{Tr}\struck{\ill{}} \leq t$. \note{trois lignes biffées, dont un « det » et un « $\Delta B$ »} $B \geq 0$. \struck{\uncertain{On} \uncertain{écrit} \ill{} \uncertain{immédiat}}
\[
  \Bigl(\Delta B = \exp\int \log\varphi_B(t) = \exp\int_0^{t} \log\varphi_B(s) = \Delta_t(B)\Bigr)
\]
\note{le second membre médian est raturé} . \uncertain{et} \uncertain{d'autre} \uncertain{part}, \ill{} $\varphi_{|1+A|}(t) \leq \varphi_{1+|A|}(t) = 1 + \varphi_{|A|}(t)$, \ill{} \ill{}
\[
  \Delta_{\struck{t}}(1 + B) \leq \Delta(1 + |B|)
\]
\note{la page s'arrête ici ; la page 36 est blanche}

\section{Préliminaires sur les opérateurs compacts d'un espace de Hilbert}

\note{Titre de l'éditeur, formé sur le titre de sa main « Préliminaires » (page 37, écrit sur « Rappels et notations » biffé). Les pages 37 à 40 sont un autre ensemble : feuillets jaunis, deux encres et un crayon rouge qui encadre les formules essentielles. Le sujet est celui des opérateurs compacts $u : E \to F$ entre espaces de Hilbert, de leurs valeurs singulières $\rho_i(u)$ et des classes $L^{p}(E,F)$ ; il ne se rattache aux pages précédentes que par la matière. Une pagination de sa main manque ; la page 39 porte « II », ce qui laisse penser que ces quatre feuillets sont ceux qui restent d'une suite plus longue.}

\page{37}
\struck{\uncertain{Dans} \uncertain{l'inégalité} (3), \uncertain{faire} \ill{} \uncertain{tendre} $p$ \uncertain{vers} $+\infty$ \uncertain{et} $q$ \uncertain{vers} $0$, \uncertain{alors} $s = \dfrac{1}{\frac{1}{p} + \frac{1}{q}}$ \struck{\ill{}} \uncertain{tend} \uncertain{vers} $0$, \uncertain{et} \ill{} \uncertain{obtient} : \uncertain{la} \uncertain{limite} \ill{} $\rho_i(uv)$} \note{en tête de page, dans un cadre tracé au trait, un bloc de trois lignes annulé par une grande boucle ; il renvoie à une inégalité (3) qui est celle de la page 38, avec des exposants $p$, $q$}

\struck{\uncertain{Rappels} \uncertain{et} \uncertain{notations}} \emph{Préliminaires}

\uncertain{L'espace} $\bar{E}$ : \uncertain{si} $E$ \uncertain{est} \uncertain{un} \uncertain{hilbert}, $\bar{E}$ \uncertain{est} \uncertain{un} Hilbert, \uncertain{et} $E' = \bar{E}$ (*) \note{la formule est encadrée au crayon rouge}. \uncertain{Propriétés} \uncertain{de} \uncertain{convergence} \ill{} : \uncertain{compacts} : \uncertain{les} \uncertain{opérateurs} \uncertain{compacts} \uncertain{sont} \uncertain{des} \uncertain{limites} \uncertain{d'op.} \uncertain{de} \uncertain{rang} \uncertain{fini}. \uncertain{Bicontinuité} \uncertain{de} $E' \mathbin{\hat{\otimes}} F \to L(E, F)$. \uncertain{Adjoint} \uncertain{de} $u : E \to F$ \ill{} $u^{*} : F \to E$, \uncertain{propriétés} \uncertain{formelles}, \uncertain{dont} \uncertain{la} \uncertain{propriété} \uncertain{spéciale}
\[
  \text{(*)}\qquad \|u^{*}u\| = \|u\|^{2}
\]
\note{la formule est cernée d'un trait qui la relie à la ligne suivante ; le numéro (*) est celui de la page, répété}

\uncertain{Décomposition} \uncertain{spectrale} \uncertain{d'un} \uncertain{opérateur} \uncertain{hermitien} \uncertain{compact} : $u = \sum \rho_i\, \bar{e}_i \otimes e_i$, \uncertain{où} $\lambda_i \to 0$, $\lambda_i = \bar{\lambda}_i$. \ill{} $u \geq 0 \Longleftrightarrow \lambda_i \geq 0$ \ill{}. \note{la somme est écrite avec des $\rho_i$ et le commentaire avec des $\lambda_i$ ; on transcrit ce qui est écrit} \emph{Corollaire} \uncertain{Valeur} \uncertain{d'une} \uncertain{fonction} \uncertain{positive} \ill{} \uncertain{compact} \ill{}.

\struck{\uncertain{Décomposition} \uncertain{canonique}} \uncertain{Décomposition} \uncertain{polaire}. \uncertain{Soit} $u : E \to F$ \ill{} \uncertain{on} \uncertain{introduit} \note{« introduit » est ajouté au-dessus de la ligne} $\sqrt{u^{*}u}$ \ill{} \uncertain{positif} \uncertain{dans} $E$, \uncertain{et}
\[
  \|\,|u|.x\,\| = \|ux\| \ill{}, \quad \text{\uncertain{d'où}}\quad ux = U|u|x,
\]
\uncertain{où} $U$ : \uncertain{isométrie} \uncertain{définie} \ill{} \uncertain{sur} $|u|(E)$, \uncertain{et} \uncertain{prolongée} \uncertain{par} \ill{} \uncertain{en} \ill{}. \uncertain{De} \uncertain{même} \ill{} \uncertain{une} \uncertain{isométrie} \uncertain{partielle} \uncertain{de} \struck{\ill{}} \ill{} $|u|(E) \subset E$. \uncertain{Et} \ill{} $=$ \ill{}. $|u| = Vu$, \uncertain{où} $V : F \to E$ \uncertain{est} \uncertain{une} \ill{} \ill{}, \uncertain{l'adjoint} \ill{} $u(E) \subset F$. \uncertain{D'où}
\[
  \text{(***)}\qquad u = U|u| \qquad |u| = Vu \qquad (\|U\| \leq 1,\ \|V\| \leq 1)
\]
\note{cerné d'un trait} \uncertain{En} \uncertain{particulier} $\|u\| = \|\,|u|\,\|$ \note{encadré au crayon rouge} \uncertain{réunis} \ill{} $\|u^{*}u\| = \|u\|^{2}$.

\marginal{en marge gauche, à hauteur de la décomposition spectrale de $|u|$ : « Rq »}

\uncertain{Décomposition} \uncertain{spectrale} \uncertain{lemme} $|u| = \sum \rho_i\, \bar{e}_i \otimes e_i$ (\uncertain{les} $\rho_i > 0$) \ill{} \uncertain{avec}
\[
  \text{(****)}\qquad u = \sum \rho_i\, \bar{e}_i \otimes f_i, \qquad \text{\uncertain{où} } f_i = Ue_i,
\]
\note{cerné d'un trait ; le $\bar{e}_i$ porte un trait de plus, raturé} \uncertain{les} \ill{} \uncertain{orthonormés} \ill{} \uncertain{avec} \ill{} ; \uncertain{Réciproquement} \ill{} \struck{$u^{*} = \sum \rho_i\, \bar{f}_i \otimes$} \ill{} \uncertain{décomposition}, \uncertain{on} \uncertain{tire} $|u| = \sum \rho_i\, \bar{e}_i \otimes e_i$, \ill{} \uncertain{les} \uncertain{valeurs} \uncertain{propres} \uncertain{de} $|u|$, \uncertain{et} \uncertain{la} \uncertain{propriété} \uncertain{correspondante} \ill{}. \uncertain{On} \uncertain{tire}
\[
  u^{*} = \sum \rho_i\, \bar{f}_i \otimes e_i, \quad \text{\uncertain{donc}}\ \ill{}
\]
\uncertain{valeurs} \ill{} \uncertain{propres} \ill{} $|u^{*}|$ \ill{} \uncertain{les} \uncertain{mêmes} \ill{} $|u^{*}| = \sum \rho_i\, \bar{f}_i \otimes f_i$ \ill{}, \uncertain{et} \ill{} $u$ \ill{}
\note{la dernière ligne, en bas du feuillet, est perdue dans une déchirure}

\page{38}
\struck{\uncertain{Rappelons} \uncertain{la} \uncertain{caractérisation}}

\emph{Notations} $(\lambda_i(u))$ \uncertain{si} $u \in L_0(E)$ : \uncertain{c'est} \uncertain{la} \uncertain{suite} \uncertain{des} \struck{\ill{}} \uncertain{valeurs} \uncertain{propres} \uncertain{de} $u$, \uncertain{rangées} \uncertain{par} \uncertain{valeurs} \uncertain{absolues} \uncertain{décroissantes}, \uncertain{avec} \uncertain{répétition} \ill{} (\uncertain{si} \uncertain{le} \uncertain{nombre} \uncertain{des} \ill{} \uncertain{est} \uncertain{fini}, \uncertain{soit} $n$, \note{« soit $n$ » ajouté au-dessus} \uncertain{il} \uncertain{y} \uncertain{a} \uncertain{lieu} \uncertain{d'attacher} \ill{} \note{en interligne : « (\uncertain{si} \uncertain{le} \uncertain{rang} \uncertain{est} \uncertain{fini} \ill{} \uncertain{nulle} \ill{}) »} \ill{} \uncertain{à} $\lambda_i(u)$ \uncertain{une} \ill{}. \uncertain{De} \uncertain{plus}
\[
  \rho_n(u) = \lambda_n(|u|) \qquad \text{\uncertain{pour} } u \in L_0(E, F).
\]
\note{encadré au crayon rouge, ici et pour les quatre autres formules encadrées de cette page ; $L_0$ désigne les opérateurs compacts} \uncertain{Donc} $\rho_n(u)$ \uncertain{est} \uncertain{une} \uncertain{suite} \uncertain{décroissante} \uncertain{à} \uncertain{valeurs} \uncertain{positives}. \struck{\ill{} \uncertain{les} \ill{}} \uncertain{Lemme} :
\[
  |\lambda_1(u)| \leq \rho_1(u) = \|u\| \qquad\qquad \rho_n(u) = \rho_n(u^{*})
\]
\note{les deux formules sont encadrées séparément ; l'indice de la première se lit $1$}

\uncertain{Si} $h$ \uncertain{est} \uncertain{hermitien} \uncertain{positif}, \uncertain{on} \uncertain{a} :
\[
  \rho_n(h) = \operatorname{Inf}_{E_{n-1}}\ \struck{\sup_{x \in F}\|hx\|}\ \sup_{\substack{x \in E_{n-1}^{\perp}\\ \|x\| \leq 1}} (hx, x)
\]
\note{cerné d'un trait ; le sup biffé est écrit en surcharge ; l'exposant du domaine de $x$ est une lecture} \struck{$h_F$ \ill{}} \uncertain{où} $E_{n-1}$ \uncertain{parcourt} \uncertain{les} \uncertain{variétés} \ill{} \uncertain{de} \uncertain{dimension} $n-1$ \note{« parcourt » est ajouté au-dessus, et « les » à sa suite} (\uncertain{caractérisation} \ill{} $(hx, y)$ \uncertain{forme} \uncertain{hermitienne} \ill{} \uncertain{positive} \ill{}).

\marginal{en marge gauche, au crayon rouge, à hauteur de la formule pour $\rho_n(h)$ : « Rq »}

\uncertain{On} \uncertain{en} \uncertain{conclut} \ill{}
\[
  \rho_n(u)^{2} = \rho_n(u^{*}u) = \operatorname{Inf}_{F_n}\ \sup_{\substack{x \in F_n\\ \|x\| \leq 1}} \|ux\|^{2}, \quad \text{\uncertain{d'où} \uncertain{aussi}}
\]
\struck{$\rho_n(u) = \operatorname{Inf} \sup_{x \in F_n} \|ux\|$} \struck{$\rho_n(u^{*}u) = \operatorname{Inf} \sup_{x \in F_n} \|u^{*}ux\|$}
\[
  = \operatorname{Inf}_{E_n}\ \sup_{\substack{x \in E_{n-1}\\ \|x\| \leq 1}} \|hx\|^{2} \qquad (h \geq 0) \struck{\ill{}}
\]
\ill{} \uncertain{avec} $h = |u|$
\[
  \rho_n(u) = \operatorname{Inf}_{F_n}\ \sup_{\substack{x \in E_{n-1}\\ \|x\| \leq 1}} \|ux\|
\]
\note{cerné d'un trait ; les indices des variétés oscillent entre $E_{n-1}$, $E_n$ et $F_n$ d'une ligne à l'autre, on transcrit ce qui est écrit} \uncertain{On} \uncertain{en} \uncertain{conclut} \uncertain{aussitôt} :
\[
  \left\{
  \begin{array}{l}
    \rho_n(uv) \leq \|u\|\,\rho_n(v) \\
    \rho_n(uv) \leq \|v\|\,\rho_n(u)
  \end{array}\right.
\]
\note{encadré au crayon rouge} \uncertain{Et} \uncertain{de} \uncertain{même} : $\rho_n(h)$ \uncertain{est} \uncertain{fonction} \uncertain{croissante} \uncertain{de} \uncertain{l'hermitien} \uncertain{positif} $h$. \uncertain{Caractérisation} \ill{} $u \in L_0(E, F)$ \struck{\ill{}} \note{une ligne de formule entièrement gribouillée} \struck{\uncertain{On} \uncertain{a} \ill{} \uncertain{Inégalités}} \uncertain{Les} \uncertain{inégalités} \uncertain{de} \uncertain{Ky}~\uncertain{Fan} :
\[
  \left\{
  \begin{array}{l}
    \rho_{m+n-1}(AB) \leq \rho_m(A)\,\rho_n(B) \\
    \rho_{m+n-1}(A+B) \leq \rho_m(A) + \rho_n(B)
  \end{array}\right.
\]
\note{encadré au trait ; les indices se lisent $m+n-1$ ; en marge gauche, au crayon rouge, un point d'interrogation}

\page{39}
\note{Le feuillet porte « II » en tête ; il est le seul des quatre à porter un numéro. Sa moitié médiane, de la proposition à « $u = \sum \rho_i \bar{e}_i \otimes f_i$ », est annulée par trois grands zigzags qui la traversent de haut en bas ; on transcrit ce qui se lit sous les traits.}

\emph{II. Caractérisation des op. de Fredholm} : $u \in L^{1}(E, F) \Longleftrightarrow \sum \rho_i(u) < +\infty$, \uncertain{et} \uncertain{alors} \uncertain{on} \uncertain{a}
\[
  \|u\|_1 = \sum \rho_i(u)
\]
\note{encadré au crayon rouge ; le $L^{1}$ est écrit $L^{(1)}$, exposant entre parenthèses}

\emph{Corollaire 1} : \uncertain{Si} $u$ Fredholm $\Longrightarrow |u|$ Fredholm, \uncertain{et} $\|u\|_1 = \|\,|u|\,\|_1 = \operatorname{Tr}|u|$.

\emph{Corollaire 2} : \uncertain{On} \uncertain{a} \uncertain{alors}
\[
  |\operatorname{Tr} u| \leq \sum \rho_i(u)\,.
\]

\marginal{en marge gauche, en diagonale : « Notation $E \mathbin{\hat{\otimes}} F$, $L^{(1)}(E, F)$ et $n \leq \dim E, \dim F$ » ; au-dessous, souligné : « Théorie de Dixmier »}

\emph{Proposition} \ill{} \uncertain{si} $u \in L_0(E, F)$ :
\[
  \sum_{1}^{n} \rho_i(u) = \sup_{\substack{v_n \in L(F, E)\\ \|v_n\| \leq 1\\ \operatorname{rang} v_n \leq n}} |\operatorname{Tr} v_n u|
\]
\note{cerné d'un trait ; sous le sup, « $v_n$ \uncertain{quelconque} » et une seconde ligne sont biffés ; les trois conditions sont écrites à droite, derrière une barre verticale} \uncertain{i.e.} \uncertain{la} \uncertain{borne} \uncertain{sup} \ill{} \uncertain{est} \uncertain{atteinte} \ill{} \uncertain{lorsque} \uncertain{les} $v_n$ \uncertain{sont} \ill{} $v_n = p_n$ \ill{} \uncertain{dans} \uncertain{les} \ill{} \uncertain{de} \uncertain{dim} $E \geq n$. \uncertain{En} \uncertain{effet}, \uncertain{on} \uncertain{a} \ill{} \uncertain{soit} \ill{} $v_n u$ \ill{} \uncertain{avec} \ill{} \uncertain{hermitien} \ill{} \uncertain{si} \ill{}, \uncertain{alors}
\[
  \operatorname{Tr} p_n v_n = \operatorname{Tr} p_n v_n = \operatorname{Tr} p_n v_n p_n,
\]
\note{ainsi sur la page, les deux premiers membres identiques} \ill{} $p_n v_n p_n = p_n(E)$, \ill{} $|\operatorname{Tr} p_n u| = |\operatorname{Tr} v| \leq \struck{\ill{}} \sum_{1}^{n} \rho_i(v)$, \ill{} $\rho_i(v) = \rho_i(p_n u p_n) \leq \rho_i(u)$, \uncertain{d'où} $|\operatorname{Tr} p_n u| \leq \sum_{1}^{n} \rho_i(u)$. \uncertain{D'autre} \uncertain{part} \ill{} $|\operatorname{Tr} p_n u| \leq \sum_{1}^{n} \rho_i(u)$, \uncertain{et} \uncertain{si} $u = \sum_{1}^{\infty} \rho_i\, \bar{e}_i \otimes f_i$, \ill{} \uncertain{alors} $\sum_{1}^{n} \rho_i(v_n u) \leq \rho_i(u)\,\|v_n\| = \rho_i(u)$, \uncertain{l'inégalité} \ill{} $|\operatorname{Tr} v_n u| \leq \sum_{1}^{n} \rho_i(u)$, \ill{} \uncertain{et} $u = p_n$ \ill{}, \uncertain{et} $\operatorname{Tr} w = \operatorname{Tr} p_n w p_n$ \ill{}, \uncertain{on} \uncertain{sait} $w'$ \ill{} $|\operatorname{Tr} w| \leq$ \ill{} \uncertain{d'où} $|\operatorname{Tr} w'| \leq \sum_{1}^{n} \rho_i(w')$ \ill{} $\leq \rho_i(w)$ \ill{} $\rho_i(w') \leq \rho_i(w)$ \ill{} \uncertain{et} \ill{} \uncertain{la} \uncertain{borne} \ill{} \uncertain{prend} $u = \sum_{i=1}^{n} \rho_i\, \bar{e}_i \otimes f_i$, \uncertain{et} \ill{}

\marginal{en marge gauche, sideways, sur sept lignes : $|\operatorname{Tr} v_n u| \leq \sum_{1}^{n} \rho_i(u)$, \ill{} \uncertain{conséquence} \ill{} ; \uncertain{d'où} \ill{} $\rho_i(u)$ \ill{} ; \uncertain{alors} \ill{} $|\operatorname{Tr} u| \leq \sum \rho_i(u)$ \ill{} \uncertain{les} \ill{} $E_n$ \ill{} ; \uncertain{de} \uncertain{rang} $\leq n$ \ill{} $u(E)$, \uncertain{donc} ; \ill{} $|\operatorname{Tr} v_n u| \leq \sum_{1}^{n} \rho_i(u)$ \ill{} ; $v = v^{*}$ \ill{} $u(E)$ \ill{} ; \uncertain{atteinte} \ill{}}

\uncertain{prend} $v_n = \sum_{j=1}^{n} \bar{f}_j \otimes e_j$, \ill{} $v_n u = \sum_{j \leq n} \rho_j\, \bar{e}_j \otimes e_j$ (\ill{}), \uncertain{dont} \uncertain{la} \uncertain{trace} \uncertain{est} $\sum_{1}^{n} \rho_i$ ; \uncertain{si} $u$ \uncertain{était} \uncertain{supposé} \uncertain{hermitien} ($f_i = e_i$) \uncertain{on} \uncertain{aurait} \uncertain{un} \uncertain{projecteur} \uncertain{hermitien} \ill{} \uncertain{dans} \uncertain{la} \uncertain{proposition}. \note{en marge gauche de ces lignes, « (?) »}

\emph{Corollaire 3} \uncertain{Soient} $u, v \in L_0(E, F)$, \uncertain{alors}
\[
  \sum_{1}^{n} \rho_i(u + v) \leq \sum_{1}^{n} \bigl(\rho_i(u) + \rho_i(v)\bigr)
\]
\note{encadré au crayon rouge ; le numéro du corollaire est surchargé} \uncertain{et} $\|u + v\|_1 \leq \|u\|_1 + \|v\|_1$. \uncertain{On} \uncertain{voit} \uncertain{ainsi} \uncertain{que} $\sum_{i \geq 1} \rho_i$ \uncertain{est} \uncertain{une} \emph{\uncertain{norme}} \ill{} \uncertain{sur} $L_0(E, F)$. \ill{}

\page{40}
\uncertain{Définition} \uncertain{de} $L^{p}(E, F)$ \uncertain{et} \uncertain{de} \ill{}
\[
  \left\{
  \begin{array}{l}
    \|u\|_p = \Bigl(\displaystyle\sum \rho_i(u)^{p}\Bigr)^{1/p} \\[2ex]
    S_p(u) = \displaystyle\sum \bigl(\rho_i(u)\bigr)^{p}
  \end{array}\right.
\]
\note{au-dessus du $\rho_i(u)$ de la seconde ligne, un « $T$ » ; la lettre devant l'accolade est raturée} \uncertain{Cas} \uncertain{des} \uncertain{variantes} \ill{} \uncertain{avec} $p = +\infty$, $p = 1$.
\note{la page s'arrête après trois lignes ; le reste du feuillet ne porte que le verso par transparence. Le lot s'achève ici.}

\end{document}
